\documentclass[../main.tex]{subfiles}
\begin{document}
%==========================================TIÊU ĐỀ TẬP========================================
\begin{table}[h]
    \begin{adjustwidth}{-1cm}{}
        \begin{tabular}{>{\centering\arraybackslash}p{6cm}|>{\raggedright\arraybackslash}p{12.5cm}}
            \multirow{5}{6cm}
            {\\[-40pt]
                \flaregothic\fontsize{20pt}{20pt}\selectfont \centering
                \color{eptype}HISTOIRE\color{black}\\
                \freshbot\fontsize{72pt}{72pt}\selectfont
                \color{epnum}18\\[6pt]
                \cafeteria\fontsize{20pt}{20pt}\selectfont\color{epnum}(D-18)\color{black}
                \\[18pt]
                \color{black}
            }
             &
            {\fotskip\fontsize{18pt}{18pt}\selectfont \ruby{川}{かわ}\ruby{内}{ち}の\ruby{週}{しゅう}\ruby{末}{まつ}}
            \\[6pt]
             & {\cafeteria\fontsize{18pt}{18pt}\selectfont Weekend at Kawachi} \\[4pt]
             & {\vnmsans\fontsize{18pt}{18pt}\selectfont Ngày nghỉ cuối tuần ở Kawachi} \\[8pt]
             & {\caviar{\fontsize{14pt}{14pt}\selectfont Nhân vật: \emp{kuon2} \emp{ryuuhou2} \emp{kaigou2} \emp{suzuno2} \emp{naomi2} \emp{game} \emp{kson} \emp{delu} \emp{mike1} \emp{koneko}}} \\[8pt]
        \end{tabular}
    \end{adjustwidth}
\end{table}
%=========================================TRUYỆN CHÍNH========================================
\color{black}

\txtbx[2]{
    {\color{vol} \uline{Tóm tắt:}} Được nghỉ trọn một ngày hiếm hoi, Kuon rủ cả nhà đi mua sắm rồi picnic ở công viên Gito. Mọi người chia thành ba nhóm: Suzuno, Naomi và Koneko lo thực phẩm; Kson cùng Delutaya mải ngắm mô hình ở khu điện tử rồi lạc sang khu khác; còn Kuon, Kaigou và Ryuuhou mua đồ gia dụng, trong đó Ryuuhou lén chọn thêm bộ câu cá. Sau vài trục trặc nhỏ, cả nhà gặp lại nhau ở công viên. Naomi và Koneko dần thân thiết hơn, còn Ryuuhou kể chuyện cười dở đến mức Kuon phải cứu nguy. Khi cả nhà về nhà, Mikeneko vẫn đứng nép ở cầu thang; Kuon chỉ nhẹ nhàng nói rằng bữa tối có phần của cô.
}

\pov{Hikawa Kuon}

\lang{Etsunan}

Một ngày cuối tháng Ba năm 2022, sau gần một tuần mưa phùn rả rích kéo dài, Kawachi cuối cùng cũng mở ra một buổi sáng trong trẻo lạ thường. Ánh nắng đầu tiên xuyên qua mây xám, rọi nhẹ vào cửa sổ phòng tôi, đánh thức tôi dậy mà không cần bất kỳ âm thanh nào. Không có chuông báo thức inh ỏi điểm đúng 6 giờ sáng như mọi ngày làm việc, không có cuộc gọi khẩn cấp từ văn phòng đòi hỏi phải xử lý hồ liệu gấp, cũng thậm chí không có tiếng Game nhỏ giọng nhắc nhở lịch trình hôm nay.

Chính sự yên tĩnh hoàn toàn ấy, thứ mà tôi hiếm khi được trải nghiệm, lại khiến tôi mở mắt sớm hơn cả khi có chuông. Tôi nằm yên trên giường, mắt nhìn chằm chằm vào vết nứt nhỏ trên trần nhà, đầu óc cố gắng lục lọi xem mình đã quên mất việc gì quan trọng đến thế mà không có bất kỳ công việc nào chờ đợi. Không có cuộc họp dài ba tiếng với đối tác, không có báo cáo tài chính cuối kỳ phải nộp, cũng không có chuyến đi công tác kéo dài cả ngày đến vùng xa xôi. Trong lúc vẫn còn hoang mang, tôi quay đầu nhìn chiếc đồng hồ treo trên tủ đầu giường, kim giờ vừa chỉ qua 7 giờ sáng.

``Kuon-user, cậu đã lật lịch điện tử, kiểm tra lịch trên bàn làm việc và ngó đồng hồ treo tường ba lần rồi đấy,'' giọng đều đều của Game vang lên từ loa nhỏ trong phòng, làm tôi giật mình.

``Tôi chỉ muốn chắc rằng ngày nghỉ này không phải là giấc mơ đâu. Quãng thời gian qua quá bận rộn đến mức tôi gần như quên mất cảm giác được ngủ đủ giấc mà không có lo âu.''

``Tất cả các phương tiện ghi nhận thời gian mà cậu sở hữu đều đồng nhất: hôm nay là thứ Bảy, không có bất kỳ lịch trình nào được lập lên cả.''

``Vậy thì tốt quá,'' tôi thở dài một hơi nhẹ nhõm, cuối cùng cũng tin rằng ngày nghỉ hiếm hoi này là thật.

Tôi ngồi dậy trên giường, duỗi đôi vai căng cứng sau cả tuần làm việc miệt mài, rồi bước chân trần xuống sàn nhà gỗ mát lạnh. Vừa mới mặc xong chiếc áo phông đơn giản, từ hành lang phía ngoài đã vọng vào tiếng chân chạy vội vã lên cầu thang, tiếng dép lách cách quen thuộc mà tôi nhận ra ngay là của ai.

Chẳng mấy chốc, Ryuuhou đã xuất hiện ở khung cửa phòng tôi, tay cầm chặt một tờ giấy được gấp cẩn thận làm tư, mặt đầy vẻ háo hức không giấu được: ``Anh Kuon, hôm nay anh có kế hoạch gì đặc biệt không?''

``Anh định đi xuống ăn sáng trước đã,'' tôi cười, kéo ghế ở bàn làm việc lại để ngồi xuống.

``Thế thì sau bữa sáng thì sao? Anh có định làm gì không?'' em vẫn đứng nguyên ở cửa, mắt lấp lánh.

Tôi liếc nhìn tờ giấy gấp tư trong tay em, nét cười càng rộng thêm. ``Cái đó là danh sách những thứ em muốn mua phải không? Đừng giấu nữa.''

Ryuuhou vội vàng giấu tờ giấy sau lưng, mặt hơi đỏ lên vì bị phát hiện: ``Em không có đâu\dots\ Em chỉ muốn nghe xem anh có ý tưởng nào hay hơn để đi chơi không thôi.''

``Để anh đoán thêm nhé: em định mua những thứ đó rồi nhờ anh trả tiền, đúng không?'' tôi nghiêng đầu, giả vờ suy nghĩ nghiêm túc.

``Em đâu có nói thế!''

``Em cũng đâu có phủ nhận là sai đâu,'' tôi đáp, nhún vai một cách thong thả. ``Được thôi, anh sẽ xem danh sách của em sau khi chúng ta ăn sáng xong. Không sao đâu.''

Ryuuhou lập tức đỏ bừng cả mặt, gật đầu liên tục rồi quay người chạy vội xuống cầu thang trước, để lại tiếng cười nhỏ vang vọng. Tôi đứng dậy, chỉnh lại áo quần rồi cũng bước theo em xuống phòng khách ở tầng dưới.

Vừa đến lan can, tôi đã thấy Kaigou đứng bên cửa sổ lớn, tay cầm điện thoại xem dự báo thời tiết một cách tập trung.

Suzuno thì đang bận rộn trong bếp, từ đó vọng ra tiếng xào nấu nhỏ nhẹ và mùi thơm của bánh mì nướng lan tỏa khắp phòng. Naomi ngồi trên tấm thảm len giữa phòng, đôi tay nhỏ bé đang xếp những chiếc lá ép khô thành hình ngôi sao đầy màu sắc, miệng còn lẩm bẩm đếm số lá một cách vui vẻ.

Kson và Delutaya thì vẫn chưa thấy bóng dáng đâu, chắc hai người vẫn đang chìm trong giấc ngủ ngon ở tầng của mình. Còn Mikeneko và Koneko, hai chị em nhà Mèo vẫn ở tầng tám, một nơi họ thường ít khi xuống trừ khi có việc cực kỳ quan trọng, thỉnh thoảng lắm mới chịu ngóc đầu ra gặp mọi người.

Tôi dựa lan can, nhìn ngắm khung cảnh quen thuộc mà ấm áp ấy, đột nhiên chợt nhận ra rằng chúng tôi đã quá lâu rồi chưa từng có dịp cùng nhau ra ngoài, mà không phải vì công việc bận rộn hay một cuộc hẹn cụ thể nào. Tất cả những lần ra ngoài trước đây đều gắn liền với trách nhiệm, với những việc phải hoàn thành, chứ chưa bao giờ là một buổi đi chơi đơn thuần.

``Này mọi người,'' tôi gọi lớn xuống phòng khách, thu hút sự chú ý của mọi người,``có ai muốn chúng ta cùng đi chơi cả ngày hôm nay không?''

Kaigou lập tức quay đầu từ cửa sổ, mặt hơi ngạc nhiên: ``Đi đâu thế? Tự dưng cao hứng à?''

``Ra khu trung tâm Kawachi thôi. Chúng ta có thể mua vài thứ đồ dùng cần thiết cho nhà, sau đó cùng nhau ăn trưa ngoài công viên, rồi về nhà trước khi trời tối. Kiểu như một ngày nghỉ đúng nghĩa, không có công việc, không có lịch trình gò bó gì cả,'' tôi nói, bước xuống bậc thang cuối cùng.

Vừa mới dứt lời, Ryuuhou đã ló đầu ra từ phòng ăn, mặt đầy phấn khích, giọng nói lanh lảnh vang lên ngay lập tức: ``Em đồng ý 100\%! Em đi!''

``Em còn chưa biết chúng ta sẽ làm gì cụ thể mà đã đồng ý rồi?'' tôi nhìn em.

``Miễn là được ra ngoài đi chơi với cả nhà, em đi đâu cũng được hết!'' em tôi nhảy cẫng lên, không giấu được niềm vui.

Kaigou khoanh tay trước ngực, nhưng khóe miệng cũng cong lên một chút. ``Thế thì phải có danh sách mua đồ rõ ràng ngay từ bây giờ. Chị không muốn tình trạng đi siêu thị rồi mới đứng giữa lối đi nhớ ra mình quên mua thứ quan trọng, lại phải quay lại đi tìm nữa.''

Suzuno lúc này thò đầu ra từ cửa bếp, mặt cũng đầy hào hứng, tay còn cầm chiếc thìa nhỏ. ``Em sẽ viết danh sách thực phẩm đầy đủ ạ! Em đã nhớ hết những thứ chúng ta sắp hết rồi!''

Naomi cũng ngừng xếp lá, giơ lên chiếc lá hình trái tim mình vừa xếp xong, mặt rạng rỡ: ``Naomi muốn đi công viên quá! Naomi muốn chạy nhảy trên cỏ!''

Tôi gật đầu, cảm thấy lòng ấm áp trước sự hào hứng của mọi người. ``Vậy thì quyết định như thế. Chúng ta sẽ chia thành các nhóm nhỏ để mua sắm cho nhanh, rồi tất cả gặp nhau ở công viên đúng mười hai giờ trưa.''

Game lúc này phát ra một tiếng bíp ngắn từ loa, giọng nói vui vẻ hơn thường ngày. ``Tôi sẽ ngay lập tức kiểm tra tất cả các tuyến xe buýt đi qua trung tâm và thời gian di chuyển chính xác, để chúng ta không phải chờ lâu ở trạm.''

``Cảm ơn ông nhiều lắm,'' tôi nói với Game, rồi quay sang thấy Kaigou đang liếc tôi với vẻ biết rõ mọi chuyện.

``Anh đừng vui quá sớm. Anh sẽ đổi ý ngay sau khi thấy danh sách mua sắm của Suzu đấy. Cô bé ấy luôn viết đủ thứ, không thiếu một món nào đâu,'' Kaigou cười, nói nhỏ với tôi.

\ct

Bữa sáng hôm ấy diễn ra nhanh hơn thường lệ, không phải vì ai vội mà vì cả nhà bắt đầu bàn xem cần mua những gì. Suzuno mở cuốn sổ nhỏ, liệt kê gạo, trứng, rau xanh, sữa, nước rửa bát và giấy ăn. Kaigou hỏi số lượng từng món, còn Naomi giúp đếm những hộp sữa còn trong tủ lạnh.

Ryuuhou nhanh nhảu kéo chiếc ghế đẩu sát mép bàn, đặt tờ giấy nhám đã gấp nhiều lần của mình vào giữa bàn, ngón tay gõ nhẹ vào từng dòng viết nguệch ngoạc. ``Em cần dây đàn mới cho cây guitar ở phòng nghe nhạc, vài cục pin dự phòng loại lớn để sạc đèn pin, cuộn băng dính vải dai loại chuyên dụng, và quan trọng nhất là một bộ lưỡi câu hoàn chỉnh mới toanh.''

Kaigou dừng tay đang cầm lọ mứt dâu, mắt hướng về tờ giấy, giọng đầy nghi ngờ. ``Lưỡi câu? Em biết câu cá à?''

Em gái tôi không hề nao núng, vỗ ngực tự hào. ``Dĩ nhiên! Bộ cũ của em để trong thùng đồ dưới sân bị ẩm mốc, lưỡi câu đã gỉ sắt hết cả rồi, không dùng được nữa.''

Cô em song sinh đặt lọ mứt xuống bàn, tiếng động nhẹ làm Naomi ngẩng lên nhìn. Tôi cất tiếng mỉa mai: ``Em có bao giờ chạm vào cây cần câu đâu mà đòi đổi lưỡi mới. Cả năm qua bộ đó nằm im dưới sân, em có bao giờ thèm đi câu đâu.''

``Lần này em sẽ đi câu thật!'' Ryuuhou nhún vai, giọng đầy hào hứng. ``Đâu biết ngày nào trời đẹp, mình cùng cả nhà đi bờ sông, lúc đó có bộ đồ sẵn thì lại không phải vội vàng mua gấp.''

Tôi đang nhai dở miếng bánh mì ốp la, nghe vậy phải chen ngang: ``Em biết cách thả lưỡi câu không? Biết cách buộc dây cước vào lưỡi không? Đừng nói là mua về để đập đi làm mẩu nhôm đốt lửa đấy.''

Ryuuhou ngẩng cao đầu, mắt sáng rực: ``Em sẽ học! You*ube đầy đủ các video hướng dẫn, mất mười lăm phút là em biết làm hết.''

Tôi lắc đầu, chỉ vào danh sách đồ dùng gia đình Suzuno đang viết: ``Đồ cần thiết cho nhà cửa trước đi, rồi mới đến mấy sở thích cá nhân của em. Bộ đồ câu cá không phải thứ cấp thiết gì trong tuần này đâu.''

Ryuuhou lập tức tỉnh bơ, chống tay lên bàn, lý lẽ đầy đủ: ``Bộ câu cá cũng phục vụ gia đình mà! Nếu em câu được mấy con cá tươi từ sông, nhà mình có thêm đồ ăn, không phải đi siêu thị mua cá đông lạnh nữa! Anh nghĩ đi, cá tự câu tươi ngon hơn hẳn!''

Kaigou nghe xong bật cười lớn, tiếng cười vang cả phòng bếp, làm Suzuno phải giơ tay che miệng để không làm đổ ly sữa: ``Được thôi, nếu em thực sự nghiêm túc câu được con cá nào về còn sống nhăn, chị sẽ đứng bếp làm sashimi cá tươi cho cả nhà ăn, em muốn ăn gì cũng làm hết.''

Mắt cô em gái lập tức sáng bừng như có đèn điện bật lên, tay vỗ bàn liên tục. ``Chị nói thật nhé! Không được đổi ý đâu đấy! Em sẽ ghi vào sổ nhớ, anh Kuon làm chứng cho em!''

``Chị mới nói là nếu em câu được cá về thôi,'' Kaigou nhấn mạnh, chỉ ngón tay vào em. ``Chưa chắc em có thể kéo nổi một con cá nhỏ lên bờ, đừng nói gì đến sashimi.''

Suzuno bên kia vẫn ghi nhanh vào cuốn sổ, vai run lên vì cố nhịn cười, bút bi lướt nhanh trên giấy để ghi hết các món cần mua. Naomi lúc này ngừng nhấm nháp miếng dâu, ngẩng lên nhìn tôi, giọng nhỏ nhẹ. ``Naomi muốn mua một quả bóng nhỏ nhựa, màu xanh lá, để chơi đá bóng hay chuyền bóng với Koneko ở công viên ạ.''

Tôi lau miệng bằng khăn giấy, nhìn em nhỏ: ``Em có hỏi em ấy là có muốn chơi không? Đừng mua về rồi con bé không thích, lại phí tiền.''

Cô em út nhanh nhảu lắc đầu, mái tóc đung đưa theo: ``Naomi sẽ lên tầng tám hỏi cậu ấy sau khi ăn sáng xong ạ. Cậu ấy hay thích ngồi đọc sách ở hiên cửa, Naomi sẽ hỏi nhẹ nhàng thôi.''

Kaigou lúc này đặt điện thoại xuống bàn, màn hình vẫn hiển thị trang dự báo thời tiết, các biểu tượng nắng vàng trải dài cả ngày. ``Dự báo thời tiết nói hôm nay trời đẹp suốt ngày, đến tận chiều tối mới có mưa nhỏ. Công viên Gito (義都) lại có con hồ, gió thổi mạnh hơn ở khu nhà mình, mọi người nhớ mang theo chiếc áo khoác mỏng, đừng để bị lạnh.''

Chúng tôi ngồi quanh bàn bàn bạc thêm mười phút, cuối cùng thống nhất là đi xe buýt công cộng thay vì lái xe nhà ra khỏi phố. Lái xe vào trung tâm cuối tuần dễ bị kẹt xe, lại khó tìm chỗ đậu. 

Game nhanh chóng truy xuất dữ liệu tuyến xe, giọng đều đều vang từ loa nhỏ trên tường. ``Tôi đã tìm được tuyến số 12, điểm dừng cuối cách cổng công viên Gito đúng bốn trăm mét, đi bộ năm phút là tới. Xe sẽ đến trạm ở đường lớn vào lúc tám giờ ba mươi.''

Suzuno liền viết thêm dòng giờ xe chạy vào cuối trang danh sách, bút màu đỏ gạch chân để không bị quên.

\ct

Sau khi ăn sáng xong xuôi, ai cũng quay về phòng chuẩn bị đồ đạc, tôi bảo Ryuuhou lên tầng tám gọi hai chị em nhà Mèo chuẩn bị. ``Đừng quên nhắc cả hai đứa mang theo áo khoác mỏng đấy, công viên có gió mạnh,'' tôi dặn thêm trước khi em bước lên cầu thang. 

Ryuuhou ngoái đầu gật đầu: ``Yên tâm anh, em sẽ nhớ mà!''

\pov{-}

\lang{Nhật}

Ryuuhou trèo cầu thang lên tầng tám, gõ nhẹ vào cửa phòng của Mikeneko và Koneko. Tiếng gõ vừa dứt, một giọng nhỏ từ trong vọng ra: ``Ai đấy ạ?''

Cô nàng đáp lại: ``Là em đây, Mikeneko-san, Koneko-chan. Onii bảo cả nhà chuẩn bị đi chơi thôi! Hai người có muốn đi không ạ?'' 

Khi cửa mở ra, bé Mèo con mỉm cười với Ryuuhou. ``Naomi đã bảo em là có bóng màu xanh lá để chơi chuyền ở công viên đó, em đã muốn đi từ lâu rồi!'' Koneko nói nhỏ, tay vân vê góc áo. 

Cô chị thì lắc đầu, dựa vào khung cửa, giọng nhỏ nhẹ nói: ``Chị mới mượn được cuốn tiểu thuyết ở thư viện tuần trước, còn chưa đọc hết một nửa. Ở nhà tự do hơn, cứ để bọn em đi chơi thôi.''

``Không sao đâu chị, nếu chị đổi ý thì nhắn tin cho em, em sẽ chạy ra cửa hàng mua thêm vé cho chị luôn, chúng ta còn kịp gặp nhau ở công viên.'' 

Naomi đứng sau lưng Ryuuhou, vừa chạy lên đến nơi, thở hổn hển, ngước nhìn Mikeneko, giọng ngoan ngoãn hỏi: ``Chị có thích ăn daikon vị mặn không ạ? Em thấy cửa hàng đó có bán gói nhỏ, em mua cho chị nhé?''

Mikeneko đứng im nghĩ một lúc, tay vắt nhẹ trên bụng, rồi lắc đầu xoa nhẹ đầu Naomi, nụ cười nhỏ hiện lên trên môi: ``Chị không cần gì đâu, em cứ đi chơi vui vẻ với Koneko là được rồi. Nhớ đừng chạy nhảy quá nhiều ở công viên, kẻo bị ngã đấy nhé,''

Koneko nhìn chị mình một lúc, ánh mắt nhẹ nhàng, rồi quay sang Naomi, nụ cười nhẹ nở trên môi. ``Tớ sẽ đi với cậu, chúng ta cùng chơi bóng ở công viên nhé.''

Cô em út nhà Hkawa lập tức mỉm cười rạng rỡ, hai má ửng hồng, tay nắm nhẹ vào góc áo Koneko. ``Vậy chúng ta cùng đi thôi, tớ sẽ dạy cậu chơi chuyền bóng, rất dễ dàng thôi.''

Mikeneko vẫn đứng dựa vào cửa phòng, tay bắt chéo trước ngực, giọng có chút cảnh báo nhưng vẫn ấm áp. ``Đừng có mua thừa đồ ăn vặt ở ngoài, rồi lại bắt chị phải ăn hết mấy cái thừa đấy, chị không chịu đâu. Lần trước cả nhà đi picnic để lại cả hộp bánh quy, chị ăn mất ba ngày mới hết.''

Naomi gật đầu liên tục, tay giơ lên như thề. ``Em sẽ nhớ ạ, không mua thừa đâu.''

``Chị không cần biết đâu nhé,'' Mikeneko nhún vai, nhưng khóe miệng đã cong lên. ``Nếu có thừa thì vẫn là chị phải ăn, chị biết mà.''

Ryuuhou nhấc chiếc túi vải nhỏ lên vai, vẫy tay với nàng Mèo hồng: ``Chị cứ ở nhà nghỉ ngơi thoải mái. Bọn em sẽ về nhà trước khi bữa tối bắt đầu, không đi chơi quá trễ đâu.''

Mikeneko gật đầu, tay vẫy nhẹ với mấy đứa. ``Chị biết rồi, các em cứ đi đi, chơi vui vẻ nhé. Đường đi cẩn thận.''

Naomi quay lại vẫy tay trước khi ra khỏi phòng. Koneko vẫy tay chậm hơn một chút. Mikeneko đáp lại bằng cái gật đầu gần như không nhận ra. Có những lời mời chỉ nên nói một lần. Nếu người ta chưa sẵn sàng, cứ để cánh cửa mở là đủ.

\ct

\pov{Hikawa Kuon}

Chặng đầu tiên của chuyến đi hôm nay là trạm xe buýt nằm cách căn nhà chung của chúng tôi khoảng tám trăm mét. Trên con đường đi xuống trạm, từng người một đều bận rộn với việc chuẩn bị cuối cùng trước khi lên xe. Suzuno vẫn luôn tỉ mỉ, bước đi mắt vẫn không rời khỏi cuốn sổ nhỏ đựng danh sách mua sắm, vừa đi vừa lướt qua từng dòng, kiểm tra lại lần cuối xem có quên mất món nào quan trọng không.

Bên cạnh em, Kaigou đang đếm lại số túi vải chúng tôi mang theo để đựng đồ, đảm bảo không có túi nào bị bỏ quên ở nhà. Còn bé Naomi thì cứ ôm chặt lấy quả bóng nhựa màu xanh lá mà em đã mang từ nhà, đôi mắt nhỏ sáng rực vì háo hức. Koneko đứng lặng bên cạnh em, không nói nhiều, chỉ đứng yên mà quan sát những hàng cây, những ngôi nhà quen thuộc trên con phố mà em ít khi có dịp dạo bước cùng cả nhà.

Riêng Ryuuhou thì đeo trên lưng một chiếc ba lô hoàn toàn trống không, khi tôi hỏi mới trả lời rằng em cần phải để trống chỗ để đựng đồ sắm được.

``\vie{Em định mua bao nhiêu thứ mà phải mang cả một chiếc ba lô to đấy?}'' Kaigou bất giác hỏi, vừa dừng lại để xếp các túi vải gọn gàng.

``\vie{Em cũng chưa biết chắc mình sẽ mua được những gì đâu, nhưng cứ chuẩn bị sẵn cho chắc thôi\~{}}'' Ryuuhou nhún vai, vẻ mặt vô cùng thản nhiên.

``\vie{Nghe em nói cứ như chúng ta đang chuẩn bị cho một cuộc sơ tán khẩn cấp chứ không phải đi mua sắm dạo cuối tuần đâu,}'' tôi tặc lưỡi cười lớn, khiến cả nhóm cùng bật cười theo.

Vừa lúc đó, hai người cuối cùng mới kịp đuổi kịp chúng tôi: Kson và Delutaya. Kson đeo cặp kính tròn quen thuộc, vai khoác chiếc áo ngắn màu đen phóng khoáng, còn Delutaya thì đeo chiếc túi chéo quen thuộc trước ngực, hai mắt vẫn còn hơi ngái ngủ vì vừa thức dậy không lâu.

``Mọi người có phải chờ lâu rồi không ạ?'' Delutaya vội vàng chạy tới, thở hổn hển hỏi.

``Cũng không lâu lắm. Vừa đủ lâu để Ryuuhou tuyên bố rằng chuyến mua sắm hôm nay là một chiến dịch lớn phải chuẩn bị kỹ lưỡng,'' tôi quay sang đáp, vỗ vào chiếc ba lô rỗng của em gái.

``Em nói đúng mà, đây thực sự là một chiến dịch thật!'' Ryuuhou lập tức phản đối, vỗ vào ba lô của mình một cách tự hào.

Kson cũng tới vỗ vai Ryuuhou một cách đồng tình. ``Phải thế mới đúng! Ai tìm được cửa hàng thú vị nhất, có nhiều đồ hay nhất trong buổi đi mua sắm hôm nay sẽ được quyền chọn món tráng miệng đầu tiên!''

Ngay lập tức, Suzuno giơ tay lên như ở trong lớp học, lắc đầu nói: ``Em xin phép rút khỏi cuộc thi đó ạ. Em chỉ có nhiệm vụ mua đúng những thứ đã ghi trong danh sách thôi, không dám đi lạc sang các cửa hàng khác đâu.''

Kaigou nhìn em gái nhỏ rồi nói: ``Đồng ý. Chị cũng không muốn phí tiền vào mấy món đồ mua về chỉ để trưng đâu.''

Vừa lúc đó, giọng đều đều của Game vang lên từ loa nhỏ trên điện thoại của tôi: ``Xe buýt tuyến số 12 sẽ đến trạm trong ba phút nữa, mọi người hãy chuẩn bị lên xe.''

Chúng tôi lập tức xếp thành một hàng ngay dưới mái che của trạm xe buýt, tránh nắng. Naomi đứng sát bên Koneko, hai em nhỏ cứ thì thầm về quả bóng mà mình mang theo. Ryuuhou đứng cạnh Kaigou, vẫn còn lẩm bẩm về những thứ em dự định mua. Còn Kson và Delutaya thì lại tiếp tục cuộc tranh luận nổ ra từ lúc ở nhà: rằng nên ghé cửa hàng điện tử trước hay là phải mua xong thực phẩm rồi mới đến đó.

Tôi lặng lẽ đứng một bên, quyết định không dính vào cuộc tranh luận vô duyên đó của hai người. Dù sao hôm nay cũng là ngày nghỉ hiếm hoi của cả nhà, tôi muốn mọi thứ cứ diễn ra tự nhiên một chút, không cần phải lên kế hoạch quá chặt chẽ như những ngày làm việc.

\ct

Chuyến xe buýt lăn bánh qua những con phố rợp bóng cây, đưa cả nhóm chúng tôi đến trung tâm Kawachi sau đúng hai mươi phút di chuyển. Không khí cuối tuần nơi đây náo nhiệt hơn nhiều so với những gì tôi tưởng tượng. Vừa bước xuống khỏi ga xe, dòng người đi bộ đã dày đặc, mỗi người một hướng tản ra khắp các con đường lớn nhỏ; những quán cà phê ven đường đã dàn ra hàng chục chiếc bàn ngoài hiên, khách ngồi đó cười nói rộn ràng, còn lũ trẻ con thì chạy tung tăng theo những quả bóng bay nhiều màu, tiếng cười vang vọng khắp không khí.

Naomi bật bước đứng lại, ngước cổ nhìn ngắm dãy bảng hiệu neon rực rỡ của các tòa nhà cao tầng mọc san sát nhau. ``Nhiều màu sắc quá Onii-chan,'' em thốt lên, đôi mắt to tròn sáng rực lên như bắt được những vì sao trên trời.

Koneko cũng dừng bước theo em, cùng hướng mắt nhìn lên những biển hiệu lấp lánh kia. Cậu bé hỏi nhỏ, giọng vẫn còn chút ngại ngần nhưng ấm áp: ``Cậu thích màu nào nhất trong số đó, Naomi-chan?''

Cô em út chống cằm suy nghĩ một lát, ngón tay nhỏ bé chạm vào cằm như đang cân nhắc thật kỹ câu hỏi quan trọng này. ``Màu xanh lá cây ạ,'' em trả lời dứt khoát.

``Tớ cũng thích màu đó nhất,'' Koneko mỉm cười, đôi tai mèo khẽ rung lên vui vẻ khi tìm được sở thích chung.

Tôi gọi cả nhóm lại sát nhau, rồi chia mọi người thành ba nhóm đúng như đã bàn từ sáng sớm ở nhà. Nhóm phụ trách mua thực phẩm sẽ gồm Suzuno cầm danh sách, cùng với hai em nhỏ Naomi và Koneko đi cùng để giúp đỡ. Nhóm chuyên đi khám phá khu điện tử vẫn là đôi bạn cùng cảnh ngộ Kson và Delutaya, hai người chắc chắn sẽ mất hàng tiếng đồng hồ để ngắm nghía đủ thứ mô hình và thiết bị ở đó. 

Còn lại nhóm mua đồ gia dụng sẽ có tôi, Kaigou và cô em gái háo hức Ryuuhou. Điểm hẹn của cả ba nhóm là cổng phía nam công viên Gito, đúng vào lúc mười hai giờ trưa, không ai được trễ.

Trước khi chia tay, mỗi nhóm tôi đều trao cho một tấm bản đồ giấy in sẵn tuyến đường, cùng số điện thoại của tôi để tiện liên lạc nếu có vấn đề phát sinh.

``Nếu có bất kỳ thay đổi gì về kế hoạch, hãy nhắn tin ngay vào nhóm chung nhé,'' Kaigou nhẹ nhàng nhắc nhở hai nhóm nhỏ, giọng chị lớn luôn cẩn thận đến từng chi tiết. ``Đừng bao giờ tự ý tách khỏi nhau đi đâu đó mà không báo cho ai biết, nghe chửa?''

``\textbf{Rõ ạ!}'' cả hai nhóm đồng thanh đáp lại, giọng nói đầy hào hứng như những học sinh chuẩn bị được đi tham quan.

Ryuuhou vỗ mạnh vào chiếc ba lô to lớn vẫn còn trống không trên lưng em, mặt đầy vẻ tự tin như một vị tướng chuẩn bị ra trận. ``Nhóm của em sẽ mang về thật nhiều chiến lợi phẩm, đợi xem nhé!''

``Chúng ta chỉ đi mua những đồ thực sự cần thiết thôi đấy,'' tôi nhanh chóng nhắc nhở em, đoán trước được cái đầu óc luôn muốn mua đủ thứ của cô em gái.

``Đồ cần thiết mà, nhưng với mỗi người thì định nghĩa về cái `cần thiết' lại khác nhau mà anh,'' em nhún vai đáp lại đầy lý lẽ, khiến tôi chỉ biết lắc đầu cười.

Ba nhóm cuối cùng cũng tỏa ra ba hướng khác nhau trên con phố đông đúc. Tôi đứng yên một lát, dõi mắt theo từng người cho đến khi tất cả đã khuất hẳn sau dòng người qua lại tấp nập. Rồi tôi cũng quay người, bước về phía cửa hàng gia dụng mà chúng ta đã định trước. Ngày nghỉ đầu tiên trong nhiều tháng của tôi, cuối cùng cũng chính thức bắt đầu.

\ct

\pov{Hikawa Suzuno}

Chỉ mất năm phút đi bộ từ trạm xe buýt, siêu thị thực phẩm lớn hiện ra trước mắt nhóm nhỏ của họ, với biển hiệu xanh lá cây rực rỡ treo trên lối vào. Tôi bước đi phía trước, cầm chặt cuốn sổ nhỏ đã ghi đầy đủ danh sách mua sắm từ sáng, ngón tay lướt nhẹ qua từng dòng chữ để không bỏ sót món nào. Naomi kéo theo chiếc giỏ đẩy có bánh xe cao su, đôi chân nhỏ bước vội vàng để theo kịp chị gái, thỉnh thoảng lại nhìn ngắm những kệ hàng đầy ắp đồ đạc bên hai bên lối đi. Koneko đi lặng lẽ bên cạnh em, đôi mắt thỉnh thoảng nhìn quanh với chút tò mò, vẫn còn hơi ngại ngần với không khí đông đúc của siêu thị cuối tuần.

Nhóm bắt đầu với những món chính nhất: bao lớn gạo trắng dày cộm được tôi đưa vào giỏ đầu tiên, sau đó là hộp trứng gà được xếp cẩn thận ở ngăn trên cùng, rồi những bó rau xanh tươi mát và thùng sữa tươi đủ loại hương vị. Trước khi đặt bất kỳ món nào vào giỏ, tôi đều lật hộp để kiểm tra ngày sản xuất và hạn sử dụng, đảm bảo tất cả đều còn mới và an toàn để ăn. 

Naomi thì phụ trách đếm số lượng hộp sữa, mỗi khi đặt một hộp vào giỏ lại chạm nhẹ vào nó và đọc thành tiếng số thứ tự, khiến Koneko phải bật cười nhỏ vì cách đếm ngây thơ của em. Khi tất cả rau củ đã được mua xong, Koneko tự giác xếp những túi rau nhẹ nhàng lên trên các hộp trứng và thùng sữa, cẩn thận để không làm dập bất kỳ lá rau nào hay làm vỡ trứng.

Khi đi qua quầy bánh ngọt và trái cây tươi, cô em út đột nhiên ngừng xe giỏ lại, quay sang hỏi Koneko với giọng vui vẻ: ``Koneko-chan, cậu muốn chọn món gì đặc biệt cho buổi picnic của chúng ta không?''

Bé Mèo con lưỡng lự, đôi mắt nhìn qua hàng ngàn món bánh và trái cây trước mặt, không biết phải chọn gì. Naomi liền gợi ý nhẹ nhàng: ``Có bánh mì nhân các loại, hay là trái cây tươi như quýt, táo, nho ấy. Cậu thích cái nào hơn?''

Koneko nhìn qua quầy bánh mì nóng hổi đang bày bán, rồi lại liếc qua những quả quýt vàng ươm xếp thành hàng đẹp mắt trên quầy trái cây, sau đó mới nhỏ nhẹ đáp: ``Tớ\dots\ muốn ăn bánh mì.''

Tôi lúc này cũng quay lại từ kệ hàng bên cạnh, mỉm cười hỏi em: ``Thế em thích loại nhân nào nhất? Ở đây có nhân đậu đỏ, nhân khoai môn, hay nhân kem đều có hết.''

Koneko chỉ tay vào chiếc bánh mì cuộn tròn, lớp vỏ vàng ươm phủ một lớp mè trắng nhỏ, bên trong là nhân đậu đỏ ngọt dịu mà em đã để ý từ nãy. ``Em muốn loại này ạ.''

Naomi gật đầu liên tục, thích thú: ``Tớ cũng thích loại đó! Vậy chúng ta lấy hai túi nhé, mỗi người một cái.''

Tôi cười nhẹ với sự đồng tình của hai em, rồi lấy hai túi bánh mì cuộn đậu đỏ đặt thêm vào giỏ, xếp nhẹ nhàng trên cùng của đống đồ mua sắm. Khi nhóm bước qua khu vực bán rau củ tươi sống, Naomi đột nhiên dừng bước lại, mắt nhìn chằm chằm vào một bó hành tây còn dính đầy đất ẩm ở gốc, vẻ mặt tò mò như thể vừa tìm thấy một kho báu.

Điện thoại của tôi đột nhiên reo lên, giọng đều đều của Game vang ra từ loa nhỏ: ``Naomi-user, chúng ta đang ở siêu thị để mua thực phẩm, không phải đến đây để làm khảo sát thực vật đâu đấy.''

``Naomi chỉ xem đất dính trên gốc hành có ướt không thôi ạ,'' Naomi nhanh chóng giải thích, tay chỉ nhẹ vào bó hành. ``Nếu đất còn ướt nghĩa là hành mới được nhổ lên, rất tươi.''

Game bật ra một tiếng cười nhỏ, giọng có vẻ thú vị: ``Đất ẩm vừa phải. Nhận xét của em hoàn toàn chính xác. Naomi-user rất giỏi.''

Lời phán xét của ông ta khiến Koneko bật cười lớn, tiếng cười nhỏ của em vang lên giữa tiếng ồn ào của siêu thị. Tôi đứng nhìn hai em gái nhỏ đang cười đùa cùng nhau, ánh mắt ấm áp nói: ``Hai em đã thân hơn rất nhiều so với lúc mới gặp nhỉ. Thật tốt quá.''

Nghe tôi nói vậy, Naomi liền đưa tay ra nắm chặt lấy bàn tay nhỏ của Koneko. Bé Mèo con hơi ngạc nhiên một chút, nhưng em không rút tay ra, mà ngược lại nắm nhẹ lại tay cô bạn. Hai em gái bước đi cùng nhau, vừa đẩy chiếc giỏ mua sắm qua các dãy hàng cuối cùng, vừa thì thầm nhỏ về những điều thú vị mà họ thấy ở siêu thị, tình bạn nhỏ giữa hai em càng thêm khắng khít trong buổi đi mua sắm đơn giản ấy.

\ct

\pov{-}

Ngay khi bước vào cửa hàng điện tử, Kson và Delutaya ban đầu chỉ có kế hoạch đơn giản: tìm mua một chiếc thẻ nhớ dung lượng cao và vài sợi dây nối cần thiết cho thiết bị của họ. Nhưng kế hoạch gọn gàng ấy chẳng kéo dài được năm phút, khi ngay ở lối vào, quầy trưng bày mô hình nhân vật thu hút toàn bộ sự chú ý của họ. Nàng Tam giác bất giác dừng chân giữa lối đi, và nàng Tổng trưởng cũng ngay lập tức dừng lại theo. Hai người trao đổi ánh mắt với nhau - trong ánh mắt cả hai đều có cùng sự tò mò, hào hứng - rồi cùng nhau bước nhẹ nhàng lại gần quầy kính.

``\vie{Cái mô hình này có phải là bản giới hạn không ạ?}'' Delutaya hỏi nhân viên bán hàng bạn tiếng Etsunan bập bẹ đang đứng gần đó, giọng có chút phấn khích không che giấu.

``\vie{Đúng vậy bạn ơi, mẫu này mới được nhập về sáng nay thôi, số lượng rất ít.}''

Kson cúi người thấp xuống để ngắm kỹ từng đường nét nhỏ trên mô hình, đầu lắc nhẹ ngưỡng mộ. ``Chất lượng sơn trên bộ giáp của nhân vật này làm tốt thật đấy, em nhìn cái đường viền này đi, sắc nét không hề bị lem chút nào.''

Delutaya cũng nghiêng người sang một bên, để ánh sáng từ đèn trần chiếu rõ hơn lên bề mặt mô hình. ``Màu sắc phối hợp trên bộ trang phục này giống hệt cái trong game em vừa chơi xong tuần trước, đúng là không thể tin được.''

Nghe vậy, nhân viên bán hàng lấy thêm một mẫu khác từ trong tủ kính ra, đặt cạnh nhau để hai người dễ dàng so sánh. Hai cô gái bắt đầu kiểm tra từng chi tiết, kích thước tổng thể có cân đối không, các khớp nối có lắp ráp mượt mà không, lớp sơn phủ có đồng đều trên toàn bộ mô hình. 

Kson nhanh chóng phát hiện ra một chiếc mô hình đặc biệt có thể thay đổi được tư thế, còn Delutaya lại tìm thấy bộ phụ kiện tùy chỉnh đi kèm để thay đổi trang bị cho nhân vật. Từ quầy mô hình, hai người bước sang dãy kệ trưng bày phụ kiện điện tử, rồi lại lạc bước sang khu vực trò chơi điện tử ở phía trong cửa hàng.

``Chúng ta đến đây để mua thẻ nhớ mà nhỉ?'' nàng Tam giác chợt ngẩng lên, vừa nhớ ra mục đích ban đầu của chuyến đi, và hỏi Kson với vẻ hơi ngạc nhiên về bản thân.

Kson nhìn quanh bối cảnh xa lạ xung quanh, gật đầu như vừa tỉnh giấc. ``Ừ ha. Cửa hàng này rộng và nhiều tầng thật, chúng ta đi lạc lúc nào không hay.''

``Kson-san, bây giờ mình đang ở tầng mấy rồi vậy?'' Delutaya nhìn quanh tìm bảng chỉ dẫn, nhưng chưa thấy đâu nên hỏi bạn mình.

Kson hướng mắt về tấm bảng treo ở cuối hành lang, sững người lại: ``Chết mịa. Tầng bốn rồi.''

``Quầy thẻ nhớ chúng ta cần tìm ở tầng hai mà?'' Delutaya thốt lên, hai người cùng nhìn nhau với vẻ hơi bối rối.

Họ nhanh chóng quay bước về phía thang cuốn, nhưng trên đường đi, hai người vẫn không ngừng bàn tán về những mẫu mô hình họ vừa ngắm nghía ở quầy trưng bày, liệt kê những chi tiết ấn tượng và dự định sẽ quay lại mua nếu có dịp.

\ct

\pov{Hikawa Kuon}

\lang{Etsunan}

Chúng tôi ba người bước dọc các lối đi trong cửa hàng gia dụng, bước chân chậm rãi chẳng nhanh hơn nhóm Kson và Delutaya bao nhiêu. Kaigou dừng lại ở quầy bát đĩa, cầm lên bộ hộp thủy tinh trong suốt có nắp đậy, mân mê lớp kính dày dặn trước khi quyết định cho vào giỏ hàng, thứ sẽ dùng để bảo quản thức ăn thừa trong tủ lạnh.

Còn Ryuuhou thì nhanh chóng bị thu hút bởi kệ chảo chống dính, em nhấc lên một chiếc đáy dày nhất, lắc nhẹ trong tay như đang cân nhắc trọng lượng. Tôi thì vừa đi vừa lướt qua tờ giấy ghi danh sách, kiểm tra từng món: khăn lau nhà, bóng đèn LED thay thế cho bóng hỏng ở phòng tắm, và túi rác sinh hoạt lớn đã sắp hết.

Cô em song sinh vừa xếp các hộp thủy tinh vào giỏ, vừa ngẩng lên hỏi tôi: ``Anh Kuon, anh có nghĩ mình còn thiếu món gì không mà liệt kê chưa có không?''

Tôi nghĩ ngợi một lúc, nhớ đến đống chìa khóa lộn xộn trên bàn ăn mỗi khi ai đó về nhà, bèn đáp: ``Thiếu một cái khay để đựng chìa khóa. Mấy đứa lúc nào cũng vứt chìa lung tung trên bàn ăn, hôm nào cần lại mất cả tiếng tìm.''

``Ừ nhỉ, cái kia quan trọng thật,'' cô nhìn sang kệ gỗ mỹ nghệ ở đầu lối đi. ``Thế thì lấy cái khay gỗ mun ở kia, nhỏ gọn mà vừa vặn.''

Ngay lúc ấy, Ryuuhou đặt chiếc chảo nặng trịch vừa nhấc vào giỏ hàng. Kaigou lập tức nhìn xuống, giọng có chút bất ngờ: ``Cái chảo đó đâu có trong danh sách chúng ta viết sáng nay?''

Em tôi thản nhiên vỗ vào mặt chảo: ``Chảo cũ ở nhà bị xước hết lớp chống dính rồi, thay cái mới hợp lý mà.''

``Em có bao giờ đụng đến cái chảo cũ đâu mà đòi thay mới?'' Kaigou rướn mày, tay chống hông. ``Có thấy em vào bếp lần nào đâu?''

``Thì, bây giờ em đang định học nấu ăn mà,'' Ryuuhou ngẩng cao đầu, mặt đầy tự tin. Tôi cá chắc là em ấy muốn trổ tài nấu một món gì đó ngoài mì úp ra.

``Ồ, học nấu ăn à? Từ bao giờ thế?''

``Từ hôm nay!''

Tôi bước lại nhấc thử chiếc chảo em vừa đặt vào giỏ, cảm nhận trọng lượng nặng nề của nó, rồi lắc đầu: ``Nếu em thực sự muốn học nấu ăn một cách nghiêm tũ thì mua loại nhẹ hơn thôi, đừng cứ chọn cái nặng nhất trong quầy. Nấu bữa đầu tiên mà bị mỏi tay thì bỏ dở giữa chừng cũng mệt đấy.''

Ryuuhou nghe vậy thì ngoan ngoãn nhấc chiếc chảo nhỏ hơn, nhẹ hơn ở cạnh đó, lắc thử vài lần rồi mỉm cười: ``Cái này được, vừa tay em.''

Kaigou nhìn em thử chảo rồi gật đầu chấp nhận: ``Được thôi, chị sẽ chỉ em dùng bếp, cách chỉnh lửa sao cho hợp lý. Nhưng nhớ đừng làm cháy nồi chảo, bếp ga mới thay tuần trước đấy.''

``Em đâu phải người hay làm cháy bếp đâu,'' cô em phản đối ngay, giọng có chút bất mãn.

``Chỉ là chưa ai dám cho em thử thôi,'' cô chị nhún vai đáp trả.

``Sao chị lại nỡ lòng nào đánh giá thấp em thế\~{}'' Ryuuhou chu môi lại, khiến hai em chúng tôi bật cười. Kể ra thì cũng hay, ít nhất thì em cũng cần biết một chút bếp núc, để sau này tự lo cho bản thân, chứ không thể lúc nào cũng chờ người khác nấu ăn cho mình mãi được.

Tôi bước sang kệ gỗ, lấy cái khay đựng chìa khóa mà hai người vừa nhắc đến, đặt nhẹ nhàng vào giỏ để không làm vỡ các hộp thủy tinh. Kaigou thì lại đếm lại số lượng khăn lau trong giỏ, kiểm tra xem có đủ năm cuộn như trong danh sách không. 

Lúc tôi quay đầu thì thấy Ryuuhou đã lẻn đi đâu mất, mắt nhìn sang phía cuối cửa hàng, thấy em đang lẻn bước về khu dụng cụ ngoài trời, nơi có treo tấm biển lớn ghi các đồ dùng cho hoạt động ngoài sân. 

Tôi lập tức gọi em quay lại: ``Này, Ryuuhou! Ra đấy làm cái gì?''

Em tôi ngoái đầu, giọng có chút ngại ngần: ``Em chỉ xem một chút thôi mà!''

``Em có thấy tấm biển to đùng ghi `dụng cụ câu cá' ở khu đó không?'' tôi hỏi lớn, dĩ nhiên biết em đang định làm gì.

``Thấy rồi ạ,'' Ryuuhou đáp nhỏ lại, biết mình bị phát hiện.

``Vậy thì em biết tại sao anh gọi em quay lại đúng không? Em biết là Kaigou sẽ ca thán thế nào nếu em vác về chứ?''

Ryuuhou mím môi, mặt có chút tiếc nuối nhưng vẫn ngoan ngoãn quay bước trở lại, chân chậm rãi bước về phía giỏ hàng của chúng tôi.

\ct

Khi chúng tôi dần di chuyển đến quầy thanh toán ở tầng trệt, Ryuuhou vẫn cứ dính sát vào bên cạnh tôi, ánh mắt cứ không ngừng lướt về phía khu dụng cụ ngoài trời ở cuối cửa hàng, như thể bộ câu cá kia sẽ biến mất nếu em không liên tục quan sát. Kaigou vừa lật giở hóa đơn thanh toán cho những món đồ đã được liệt kê trong danh sách, còn tôi thì lần lượt đặt những món đồ còn lại lên băng chuyền để nhân viên quét mã.

``Em vẫn còn muốn bộ câu cá đó à?'' tôi thấp giọng hỏi, nghiêng đầu nhìn em gái đang thèm thuồng nhìn xa xăm.

``Em chỉ thấy nó thực sự hữu ích cho cả nhà mà thôi,'' Ryuuhou đáp liền, vẫn không rời mắt khỏi khu vực kia.

``Em chưa từng có kinh nghiệm câu cá bao giờ đâu,'' tôi nhắc lại sự thật đơn giản mà em dường như đã quên.

``Em có thể học mà, có gì đâu mà khó,'' em quay sang nhìn tôi, mặt đầy vẻ tự tin rằng mọi thứ sẽ dễ dàng như em nghĩ.

Kaigou vừa đưa thẻ ngân hàng cho nhân viên thanh toán, vừa quay lại nghe được câu cuối của Ryuuhou, cô chị lập tức thêm vào: ``Nếu em thực sự muốn, thì em phải tự tìm hiểu mọi thứ về câu cá, và tuyệt đối không được bắt người khác phải dọn dây câu hay những thứ lộn xộn sau này, lúc này chị mới cân nhắc.''

Ryuuhou gật đầu như giã giơ, không chút do dự: ``Em hứa sẽ tự dọn dẹp tất cả, không phiền đến ai đâu.''

``Và thêm một điều nữa, không được mua bộ đắt nhất trong quầy,'' Kaigou thêm một điều khoản nữa, không cho em cơ hội né tránh.

``Em sẽ chọn bộ cơ bản nhất, đảm bảo không tốn quá nhiều tiền,'' em nhanh chóng đáp ứng, như thể đã chuẩn bị sẵn các lời hứa để thương lượng.

Tôi nhìn em với ánh mắt nghi ngờ, bật cười nhỏ: ``Em đang thương lượng với chị ấy như thể đã quyết tâm mua bộ đó từ lâu rồi, chứ chẳng phải vừa mới nghĩ đến đâu.''

``Em chỉ đang hỏi ý kiến hai anh chị thôi mà, có gì đâu,'' Ryuuhou vờn ngón tay, mặt giả vờ vô tội.

Kaigou thở dài một hơi, nhưng khóe miệng không thể nào giữ được độ cong, khóe mắt cũng nhoẻn lên: ``Thôi được rồi. Mua bộ cơ bản, nhưng tuyệt đối không được mua thêm mồi câu gì thêm đâu.''

Ryuuhou lập tức vờn ngạc nhiên, đưa tay đặt lên ngực như thể vừa nghe một điều không thể tin nổi: ``Em còn chưa từng nghĩ đến chuyện mua mồi câu đâu ạ.''

``Em vừa mới nghĩ đến rồi, đừng có giả vờ nữa,'' tôi bật cười lớn, đập nhẹ vào vai em.

Chúng tôi sau đó lại quay bước trở về quầy dụng cụ ngoài trời, Ryuuhou liền lao tới và chọn một chiếc cần câu gấp gọn dễ dàng mang theo, cuộn dây cước tiêu chuẩn và một chiếc hộp nhựa nhỏ để đựng phụ kiện. Kaigou cầm lên bộ câu kiểm tra các thông số in trên bao bì, đảm bảo chất lượng ổn, còn tôi thì mở tờ hướng dẫn sử dụng, đọc lướt qua các bước lắp ráp đơn giản. Cuối cùng, chúng tôi cũng đặt bộ dụng cụ câu cá đó vào giỏ hàng, hoàn thành ý định của con bé.

Ryuuhou mím môi cười tươi, nhìn bộ câu cá như vừa nhặt được báu vật: ``Em đã nói nó sẽ phục vụ tốt cho gia đình mà.''

``Em còn chưa câu được con cá nào đâu, đừng mừng sớm,'' tôi nhắc nhở em, nhưng cũng không thể nào không cười theo niềm vui của em.

``Nhưng bây giờ em đã có cơ hội để học rồi, đúng không nào?'' em trả lời đầy hào hứng, bước đi phía trước với vẻ phấn khích không thể che giấu.

\ct

\pov{-}

Cuối cùng cũng bước chân đúng đến tầng hai của cửa hàng điện tử, Kson và Delutaya thay vì vội vàng đi tìm quầy thẻ nhớ theo kế hoạch ban đầu, lại dừng chân bất ngờ trước dãy máy chơi game cầm tay mới ra lò. Một nhân viên bán hàng đang nhiệt tình giới thiệu những tính năng đột phá của dòng thiết bị vừa cập bến, giọng nói tràn đầy sự phấn khích lôi cuốn. Kson cầm nhẹ trên tay chiếc máy mẫu, ngón tay khẽ chạm vào bề mặt màn hình bóng loáng, còn Delutaya thì nghiêng người sát lại, ánh mắt dán chặt vào màn hình đang hiển thị hình ảnh demo sống động.

``Màn hình này độ sáng khá ổn, nhìn ngoài nắng cũng không bị chói quá. Còn chất lượng âm thanh thì sao nhỉ?'' nàng Tam giác lên tiếng hỏi, đầu vẫn chưa rời khỏi màn hình.

Nhân viên nhanh chóng đưa cho hai người một cặp tai nghe cao cấp được kết nối sẵn với máy. Hai cô gái cùng đeo lên và thử nghe một đoạn nhạc ngắn được cài sẵn, cảm nhận từng âm thanh lan tỏa.

``Âm trầm khá sâu và ổn định, nghe không bị bể tiếng,'' nàng Tổng trưởng là người đầu tiên nhận xét, gật đầu nhẹ vẻ hài lòng.

``Nhưng tiếng treble thì hơi gắt một chút, nghe lâu chắc sẽ mỏi tai,'' Delutaya bổ sung thêm, tháo tai nghe ra để đặt lại trên quầy.

Hai người rồi lại bắt đầu hỏi dồn dập về thời lượng pin, độ bền của vỏ máy khi vô tình bị rơi, và quan trọng nhất là khả năng tương thích với những bộ điều khiển cũ họ đang có ở nhà. Cuộc trò chuyện kéo dài mãi, đến mức nhân viên phải lấy thêm cả bảng thông số kỹ thuật dày cộp ra để giải thích từng chi tiết. 

Đúng lúc ấy, Kson chợt liếc nhìn đồng hồ đeo tay, mặt lập tức thay đổi. ``Delu-chan, giờ đã mười một giờ mười lăm rồi kia!''

``Cái gì? Vậy là chúng ta chỉ còn bốn mươi lăm phút để chạy đến công viên Gito đúng giờ hẹn thôi à?'' Delutaya cũng hoảng hốt nhìn đồng hồ của mình.

``Mà thẻ nhớ chúng ta cần mua thì vẫn chưa mua được đâu,'' Kson thở dài, chỉ còn biết chỉ tay về phía quầy thẻ nhớ nằm ngay phía bên kia hành lang, cách họ chưa đầy mười bước chân.

Hai người vội vàng chạy sang quầy, chỉ mất vài phút là đã hoàn tất việc thanh toán cho chiếc thẻ nhớ dung lượng cao. Nhưng vừa bước ra khỏi quầy thanh toán, mắt của nàng tóc tím lại chạm phải quầy mô hình ở góc hành lang, nơi có một chiếc phụ kiện giới hạn mà cô đã để ý từ lúc nãy. 

Cô định bước lại xem xét thêm một chút thì Delutaya nhanh chóng giơ tay ngăn trước mặt: ``Không được đâu Kson-san. Bây giờ phải đi đến điểm hẹn ngay, không thể trễ giờ được.''

``Chị chỉ nhìn thêm một chút thôi, chỉ một chút thôi mà\~{}'' Kson nũng nịu, chân vẫn chưa chịu dịch chuyển.

``Lần trước chị cũng nói chỉ nhìn một chút, mà kết quả là chúng ta đã đi lạc hết cả ba tầng liền đó!'' nàng tóc xanh lá nhắc lại tình huống không lâu trước đó, giọng đầy cảnh cáo.

``Đó là vì cửa hàng này có quá nhiều tầng thôi mà, đâu phải lỗi của chị đâu!'' Kson còn muốn biện minh, nhưng nhìn vào khuôn mặt nghiêm nghị của bạn mình thì cũng đành ngậm ngùi.

Delutaya sau đó nhẹ nhàng kéo nhẹ vạt áo của Kson, dẫn dắt bạn mình bước về phía lối ra của cửa hàng. Kson tuy ngoan ngoãn đi theo, nhưng đầu thì vẫn ngoái liên tục về phía quầy mô hình, ánh mắt đầy tiếc nuối như sắp phải xa một người thân yêu.

\ct

\pov{-}

Chúng tôi dứt điểm mọi khoản thanh toán tại cửa hàng gia dụng đúng mười một giờ ba mươi, đúng lúc điện thoại trong túi tôi reo lên liên tục, nhắn cả ba nhóm đều đã hoàn thành nhiệm vụ mua sắm của mình. Nhóm thực phẩm dẫn đầu bởi Suzuno nhắn vào nhóm chung báo đã check hết các món trong danh sách, kèm theo bức ảnh cuốn sổ nhỏ với từng ô được gạch chân rõ ràng, chẳng để sót món nào. Đằng sau tin nhắn của chị là một bức vẽ nguệch ngoạc đáng yêu của Naomi: ba chiếc bánh mì được tô màu vàng ươm, có cặp mắt tròn xoe vẽ lên mỗi chiếc trông như đang cười. Còn Koneko chỉ gửi một biểu tượng mặt cười màu cam đơn giản, đủ để mọi người hiểu mọi thứ diễn ra êm xuôi.

Trên đường ra khỏi cửa hàng, tôi xách chiếc túi chứa khăn lau và túi rác, còn Ryuuhou thì hớn hở đeo chiếc ba lô to đùng trên lưng, bên trong là bộ câu cá vừa mới mua được, em cứ sờ soạng chiếc khóa kéo liên tục như sợ nó sẽ biến mất bất cứ lúc nào.

Kaigou thì vai đeo chiếc túi lớn nhất và nặng nhất, đựng đầy những hộp thủy tinh trong suốt vừa mua cùng mấy bóng đèn LED thay thế; tất cả đều là đồ dễ vỡ, nên cô nàng phải đi chậm lại, tay giữ chặt quai túi để chúng không va đập vào nhau.

``Chị Kaigou, có muốn em cầm bớt một vài món không?'' Ryuuhou đột nhiên quay đầu hỏi, mắt nhìn thấy tôi đang khòm lưng giữ túi.

``Thôi, chả cần, em đang mang cả bộ câu cá trên lưng rồi mà,'' cô nàng lắc đầu từ chối, biết em đã có một trọng lượng không nhỏ trên vai.

``Nhưng nó nhẹ lắm à,'' em nhún vai, nhún nhảy vài bước để chứng minh lời nói của mình, chiếc ba lô không hề kêu cót két gì.

Kaigou nhìn chiếc ba lô trên lưng em, gật đầu xác nhận: ``Đúng là em chọn loại gấp gọn được nên nó nhẹ hơn nhiều so với những bộ cần câu to đùng khác. Không tốn sức mang đâu.''

``Thế mà hồi nãy mà vác cái khủng nhất ở trong đó thì có khi em thành `tể tướng Lưu Gù' rồi cũng nên,'' tôi xéo giọng mỉa mai, nhớ lại lúc ở quầy dụng cụ ngoài trời em đã suýt nữa thì nhấc chiếc cần câu dài nhất, nặng nhất trong quầy lên.

Chúng tôi bước đi dọc theo con đường dẫn đến công viên Gito, hai bên đường là hàng cây cổ thụ rợp bóng mát, che bớt đi cái nắng buổi trưa đang dần lên. Hai bên vỉa hè có nhiều quầy hàng nhỏ ven đường bán nước giải khát mát lạnh và những món ăn nhẹ thơm phức, mùi vị bay theo gió làm cho bụng tôi cũng cồn cào lên chút ít.' Kaigou lịch sự đi ở phía ngoài, sát làn đường xe cộ qua lại, để Ryuuhou đi ở phía trong, tránh cho em bị va quệt bởi những người đi bộ vội vã. Tôi đi lùi phía sau hai người, nhìn hai cô nàng trước mặt đang nhắng nhít nhau.

Được một đoạn, Ryuuhou lại quay đầu nhìn tôi, mày nhíu lại. ``Anh đi chậm quá vậy, sắp bị bỏ lại phía sau mất!'' em kêu lên, bước chân đã sẵn sàng để chạy về phía trước.

``Cứ bình tĩnh mà đi thôi, không cần phải vội quá đâu, còn tận mười lăm phút nữa cơ mà,'' tôi giải thích.

``Vậy em có thể cầm chiếc chảo trong túi anh giúp mà, để anh đỡ mang một món,'' em đề nghị ngay, như thể đã nghĩ ra một giải pháp hoàn hảo.

``Thì cái chảo đâu có nằm trong túi của anh đâu,'' tôi bật cười, lắc đầu ngán ngẩm trước sự quên nhanh của con bé. ``Em đã tự nó cất vào trong ba lô của mình rồi còn gì.''

\ct

Công viên Gito nằm bên bờ một nhánh sông nhỏ uốn lượn, với những bãi cỏ xanh mướt trải rộng cả một vùng và hàng cây anh đào đã lỡ mùa nở hoa, chỉ còn lại những cành lá non xanh tươi mát rượi. Ngay tại lối vào công viên, một tấm biển nội quy được treo trang trọng, ghi rõ quy định khách đến picnic phải giữ khoảng cách an toàn ít nhất năm mét tính từ bờ nước. 

Suzuno đã đến sớm hơn mọi người một chút để chọn trước một mảnh cỏ phẳng phiu, nằm ngay dưới bóng mát rợp của một cái cây cổ thụ, tránh được cái nắng buổi trưa đang dần trở nên gắt gỏng. Naomi và Koneko đã trải xong tấm bạt lớn lên mặt cỏ, đang bận rộn sắp xếp các túi thực phẩm thành từng hàng gọn gàng. Thấy tôi và Kaigou, Ryuuhou vừa bước tới với những bao đồ nặng trên tay, Koneko liền nhanh nhẹn đứng dậy đến phụ đỡ những chiếc túi còn lại.

``Em tìm được một vị trí cũng ngon nghẻ đấy,'' tôi nói với Suzuno, mắt nhìn quanh khoảng không gian thoáng đãng mà em ấy đã chọn.

``Em muốn đến sớm một chút để chắc chắn có được chỗ có bóng râm, tránh để mọi người phải ngồi dưới nắng nóng ạ,'' con bé cười nhẹ, đáp lại lời khen của tôi.

``Còn K và Delu thì sao? Hai đứa ấy đến đâu rồi?'' Kaigou vừa đặt chiếc túi nặng của mình xuống đất, vừa ngẩng lên hỏi.

Suzuno lướt mắt nhìn đồng hồ đeo tay, rồi đáp: ``Hai chị ấy vừa nhắn tin là đang trên đường đi bộ từ cửa hàng điện tử đến đây, sớm thôi là đến ạ.''

Lúc này Ryuuhou cũng vừa đặt chiếc ba lô của mình xuống bãi cỏ. Naomi tình cờ liếc thấy một gói nhỏ xíu bị lọt ra khỏi khe khóa kéo của ba lô, liền tò mò hỏi: ``Cái gì trong đó vậy ạ, chị Ryuuhou?''

Em tôi đột nhiên khựng lại một chút, như thể không ngờ bị phát hiện sớm thế.

Kaigou đứng kế đó, tay khoanh trước ngực với vẻ mặt như đã đoán trước mọi chuyện: ``Nói đi, em lại giấu cái gì nữa thế?''

Ryuuhou đành cười khờ, mở rộng chiếc khóa kéo và lấy ra bộ cần câu gấp gọn mà em ấy đã đòi mua từ lúc ở cửa hàng. ``Em\dots\ em mua bộ câu cá đấy.''

Naomi lập tức mở to mắt tròn xoe, giọng đầy bất ngờ: ``Vậy mình có thể câu cá ngay ở đây không ạ? Trông nó thú vị quá!''

Cô chị Hai liền chỉ tay về phía tấm biển nội quy ở lối vào mà chúng ta vừa đi qua. ``Không được đâu em, quy định công viên cấm câu cá trong toàn bộ khu vực này để bảo vệ nguồn thủy sản đấy em.''

Ryuuhou liếc nhìn tấm biển, rồi lại nhìn chằm chằm vào chiếc cần câu trong tay, giọng có chút tiếc nuối nhưng vẫn nhanh chóng gạt sang một bên: ``Em mua trước để dành cho lần sau đi đâu đó có thể câu mà, không phải định dùng ở đây đâu.''

``Thôi cũng được, ít nhất thì em cũng nhớ hỏi trước khi tự ý mở ra dùng bừa bãi,'' tôi vỗ vai em ấy, cười nhẹ.

Kaigou giúp em ấy cất lại bộ dụng cụ câu cá vào trong ba lô, rồi nói: ``Lần sau nếu thực sự muốn đi câu, chị sẽ chỉ em cách buộc dây câu cho đúng. Đừng nghĩ là mua được bộ về là có thể đi câu ngay được.''

Ryuuhou lập tức nở nụ cười đắc ý, như thể vừa giành được chiến thắng nhỏ: ``Chị đã hứa nhé! Không được rút lời đâu!''

``Chị chỉ nói là sẽ dạy em cách buộc dây thôi, chứ chưa hề nói là sẽ đi câu cùng em đâu đấy,'' cô lập tức nhấn mạnh, tránh bị con bé lật lọng.

``Em nhớ chính xác từng chữ chị nói mà, không bao giờ quên đâu,'' con bé cười hì hì, rõ ràng là đã đạt được mục đích nhỏ của mình.

Sau khi mọi thứ đã ổn thỏa, Suzuno bắt đầu mở các túi thực phẩm, lần lượt bày tất cả đồ ăn lên trên tấm bạt để chuẩn bị cho bữa picnic chung.

Trên tấm bạt trải rộng dưới bóng cây, đủ thứ thức ăn được bày biện gọn gàng: những ổ bánh mì mềm xốp, đĩa trứng cuộn vàng ươm, tô salad rau tươi mát, đĩa trái cây được cắt thành miếng vừa ăn, cùng với mấy chục hộp cơm nắm nhỏ xinh được nặn kỹ càng. Naomi chăm chỉ xếp từng chiếc đĩa giấy thành hàng ngay ngắn, còn Koneko thì tỉ mẩn gấp từng chiếc khăn ăn thành hình tam giác vuông vức, đặt cạnh mỗi chiếc đĩa. 

Tôi khiêng bình nước lớn đặt đúng chính giữa tấm bạt, đảm bảo bất kỳ ai ngồi ở đâu cũng có thể với tay chạm tới để rót. Kaigou thì ở cạnh kia, giở túi rác ra kiểm tra xem đã mở sẵn miệng túi chưa, sợ rằng mọi người ăn xong lại vất vỏ hộp bừa bãi.

``Kaigou chuẩn bị mọi thứ chu đáo thật,'' tôi vừa đặt bình nước xuống vừa nói, cảm nhận sự chu đáo của cô em gái song sinh.

``Nếu không cẩn thận thế này,'' Kaigou đáp lại, ``lát nữa lại có người vứt rác lung tung khắp nơi, dọn dẹp mệt chết đi được.''

``Naomi không bao giờ vứt rác bừa bãi đâu ạ,'' Naomi vừa xếp xong chiếc đĩa cuối cùng liền ngẩng lên khẳng định, mặt tròn xoe đầy thành thật.

``Chị biết mà, em ngoan nhất nhà. Chị đang nói đến mấy người lớn đấy, chứ đâu phải nói em đâu,'' Kaigou cười nhẹ, chỉ mấy người đang tụ tập xung quanh.

\lang{Nhật}

Vừa nói xong, Kson và Delutaya vừa thở hổn hển xuất hiện ở lối vào công viên, mỗi người tay cầm một túi nhỏ xíu. Kson còn giơ cao chiếc thẻ nhớ vừa mua được lên như đang khoe một tấm huy chương chiến thắng, miệng hét toáng lên: ``\textbf{Nhiệm vụ mua sắm hoàn thành xuất sắc!}''

Delutaya nhìn về phía chúng tôi đang ngồi dưới cây, thở phào nhẹ nhõm: ``May quá, mọi người vẫn còn ở đây, tưởng mình đến trễ mất.''

``Chị nghĩ là bọn em đi lạc đường rồi à?'' Ryuuhou từ bên cạnh cất tiếng hỏi, mái tóc hơi rũ vì chạy bộ đến.

``Không đâu. Chị chỉ sợ Kson-san lại bị cuốn hút ở quầy mô hình nhân vật nào đó nữa, quên cả thời gian đến đây,'' Delutaya giải thích, mắt còn liếc sang bạn mình đầy cảnh giác.

Kson lập tức chống nạnh, mặt đầy không phục: ``Chị tự đi được mà, đâu có bị giữ lại đâu.''

``Chị còn dám nói? Lúc ở cửa hàng chị đã quay lại quầy mô hình đó hai lần liền đó, nếu không phải em kéo đi chắc còn ở lại nữa,'' nàng Tam giác nhắc lại chuyện không lâu trước đó, khiến nàng Tổng trưởng bật cười ngượng.

``Em nhớ dai thế, chỉ một chút thôi mà cũng ghi nhớ mãi,'' Kson xoa đầu cười trừ.

Naomi thấy hai người đến liền vẫy tay rủ rê: ``Kson-oneechan, Delu-oneechan, lại đây nào, đồ ăn sắp xếp xong xuôi rồi ạ.''

Hai người nhanh nhẹn đặt túi nhỏ xuống, rửa tay sạch sẽ bằng khăn ướt rồi lần lượt ngồi vào chỗ trống trên tấm bạt. Bữa picnic chính thức bắt đầu khi mọi người cùng nhau chia sẻ từng món thức ăn. Suzuno lần lượt đưa cho mỗi người một hộp cơm nắm nhỏ, nóng hổi vừa thơm vừa ngon. Naomi liền lấy ra chiếc bánh mì nhân đậu đỏ mà mình đã chọn từ hồi ở siêu thị, Koneko cũng cầm lên một chiếc giống hệt như của bạn. Hai đứa trẻ ngồi sát cạnh nhau dưới bóng mát của cây cổ thụ, gió thổi nhẹ làm tóc hai em bay phấp phới. Naomi cẩn thận bóc lớp giấy gói bánh mì ra, không để cho một mẩu giấy nào rơi xuống cỏ.

``Cậu đã từng ăn loại bánh mì nhân đậu đỏ này bao giờ chưa?'' Naomi nhẹ nhàng hỏi bạn mình, mắt nhìn chiếc bánh trong tay đầy tò mò.

``Chưa bao giờ, cậu thì sao? Cậu có ăn thử chưa?'' Koneko lắc đầu, cũng ngắm nghía chiếc bánh trong tay mình.

``Tớ cũng chưa từng ăn, hôm nay là lần đầu tiên đấy.''

Naomi cắn một miếng nhỏ, vị ngọt của đậu đỏ hòa quyện với vị mềm của vỏ bánh lan tỏa trong miệng, em liền gật đầu phấn khích: ``Ngọt vừa phải thôi, không hề bị gắt cổ gì cả, ngon lắm.''

Koneko cũng nếm thử một miếng, mắt lập tức sáng lên, gật đầu lia lịa: ``Ừm ừm, đúng là ngon thật sự, mình chọn đúng món rồi.''

Sau đó Naomi chia cho Koneko một miếng dâu tây lớn mọng nước, Koneko cũng đáp lại bằng cách đưa cho em một miếng bánh mì to hơn. Không có ai bắt hai đứa phải nói chuyện liên tục, hai em chỉ yên bình ngồi cạnh nhau, ăn uống nhẹ nhàng, thỉnh thoảng lại chỉ tay lên trời chỉ cho nhau xem những đám mây trôi lững lờ, có đám hình con thỏ, có đám lại giống chiếc bánh kem khổng lồ. Tôi ngồi ở phía bên kia tấm bạt, nhìn cảnh tượng êm dịu đó mà lòng thấy ấm áp lạ thường. Suzuno cũng cùng tôi nhìn hai đứa trẻ, miệng nở một nụ cười dịu dàng.

``Naomi-chan hôm nay vui lắm anh ạ, từ lúc đến giờ em chưa ngừng cười đâu,'' em nhỏ giọng nói với tôi, mắt vẫn dõi theo hai em.

``Koneko-chan cũng vậy, em ấy cũng mở lòng hơn nhiều so với lúc mới gặp,'' tôi đáp lại, cảm thấy vui mừng vì hai đứa trẻ đã hòa hợp với nhau.

Kaigou ở bên cạnh cũng gật đầu đồng tình: ``Hai đứa hợp nhau hơn chị đã nghĩ ban đầu, cứ như là đã quen nhau từ lâu rồi vậy.''

Kson nghe vậy cũng giơ ngón cái lên tán thành: ``Bọn trẻ con lúc nào cũng có cách riêng để làm quen với bạn mới, không cần người lớn phải lo lắng gì cả.''

Delutaya cũng gật đầu thêm vào: ``Đúng vậy, đôi khi chỉ cần ngồi cạnh nhau im lặng thôi cũng đủ để hai người trở thành bạn thân rồi.''

Ryuuhou vừa nhai miếng trứng cuộn vừa nghe mọi người nói, liền ngẩng lên phản đối: ``Mọi người nói kiểu gì thế, em cũng chẳng tài nào hiểu được ấy.''

``Em thực sự hiểu à?'' Kaigou nghiêng đầu hỏi em gái, mặt đầy vẻ không tin.

``Em hiểu chứ! Em đã từng ngồi cạnh Onii cả tiếng đồng hồ mà không nói một câu nào đâu,'' Ryuuhou khẳng định, như thể đó là một thành tích lớn.

``Đó là vì hôm đó em giận anh, chứ đâu phải tự nhiên mà ngồi im vậy,'' tôi liền nhắc lại sự thật, khiến mọi người xung quanh bật cười.

``Dù sao thì cũng đã ngồi im cạnh nhau cả tiếng, vậy vẫn tính chứ sao,'' con bé cười khúc khích, không hề thấy ngại ngần gì.

\ct

Khi bữa trưa vừa kết thúc, những mẩu thức ăn cuối cùng cũng được dọn gọn vào hộp, Kaigou liền xòe ra một xấp thẻ giấy viết tay với đủ loại từ ngữ khác nhau, đề xuất cả nhà cùng chơi trò đoán từ để giải trí. 

Lượt đầu tiên rơi vào Ryuuhou, em rút ngẫu nhiên một chiếc thẻ và đọc lướt dòng chữ: ``cần câu''. Ngay lập tức, em ngồi thẳng người, hai tay bắt chước cầm một chiếc cần câu dài tưởng tượng, sau đó làm động tác kéo mạnh như vừa móc được con cá nặng, lắc lư hai cánh tay để mô tả cảnh vật câu kéo mệt mỏi.

Naomi, với đôi mắt tròn xoe quan sát không chớp, liền hô lên ngay trước khi kịp suy nghĩ lâu: ``Chị đang diễn tả câu cá ạ!''

Ryuuhou vỗ tay reo mừng, gật đầu lia lịa để xác nhận đáp án đúng. Đến lượt Kson tham gia, cô rút phải thẻ ghi từ ``máy chơi game'', liền ngay lập tức giơ hai tay ra trước ngực, các ngón tay bấm nhảy liên tục trên không trung như đang nhấn các nút điều khiển trên máy, thỉnh thoảng còn nghiêng người sang trái phải như đang tránh chướng ngại vật trong trò chơi.

Delutaya, người luôn ở bên cạnh nàng Tổng trưởng và hiểu rõ thói quen của em ấy, đoán ra ngay chỉ sau vài giây quan sát: ``Đó là máy chơi game cầm tay đúng không?''

``Đúng rồi! Đúng hoàn toàn!'' Kson reo lên vui sướng, thả lỏng hai tay và cười lớn vì Delutaya đã bắt kịp ý đồ diễn xuất của mình.

Đến lượt tôi, tôi rút được chiếc thẻ ghi từ ``xe buýt''. Trời ạ.

Tôi ngẫm nghĩ một chút rồi bắt đầu làm động tác: hai tay bắt chước nắm vô lăng, lắc nhẹ sang trái rồi sang phải như đang lái xe, sau đó đạp chân xuống đất như đang thắng gấp, rồi mở miệng giả vờ hô lớn: ``Ai xuống trạm tiếp theo đây!''. 

Kaigou ngồi đối diện tôi, nhíu mày quan sát từng động tác một lúc lâu, cuối cùng đưa ra dự đoán: ``Xe tải?''

``Không phải đâu.'' Tôi lắc đầu, tiếp tục lặp lại động tác lái xe.

``Xe van chở hàng?'' Kaigou thử lại câu trả lời thứ hai.

``Vẫn sai nữa.'' Tôi vừa nói vừa tiếp tục làm động tác dừng xe, mở cửa giả vờ cho khách lên xuống.

Ryuuhou ở bên cạnh không chịu nổi nữa, chen ngang vào cuộc: ``Anh đang làm gì thế kia? Mấy động tác này là cái gì vậy?''

``Là xe buýt đấy, không nhận ra à?'' tôi hơi ngạc nhiên vì không ai đoán ra.

``Anh diễn tệ quá! Diễn như vậy ai mà đoán được!'' Ryuuhou phàn nàn, cười lớn vì động tác của tôi quá khó hiểu.

``Anh đâu phải diễn viên chuyên nghiệp, đâu thường xuyên chơi trò này nên diễn không khéo thôi.'' tôi bào chữa cho mình.

``Vậy thì lần sau đừng rút thẻ nữa, làm tốn thời gian cả nhóm!'' con bé hài hước đòi loại tôi ra khỏi cuộc chơi.

Naomi nghe vậy thì bật cười to, ngã nghiêng về phía Koneko vì quá vui. Tôi giả vờ đặt tay lên ngực, mặt ra vẻ bị xúc phạm nặng nề, nhưng Kaigou đã nhanh chóng đưa một chiếc thẻ mới cho Ryuuhou để tiếp tục trò chơi. 

Lần này em rút được thẻ ghi từ ``một câu chuyện cười'', em liền ngước lên nhìn chị mình, thắc mắc: ``Từ này phải kể bằng lời ra để mọi người đoán hả chị? Cái này không thể diễn tả bằng động tác được đâu.''

Kaigou nhún vai, cười nhẹ đáp: ``Tùy em thôi, muốn diễn tả sao cũng được, miễn mọi người đoán ra là được.''

Nghe vậy, Ryuuhou liền ngồi thẳng dậy, chỉnh lại cổ áo, hắng giọng to như một người dẫn chương trình chuyên nghiệp đang đứng trên sân khấu lớn: ``Hôm nay tôi sẽ kể một câu chuyện ngắn! Có một con cá bơi bơi mãi, cuối cùng cũng vào được một thư viện ở dưới sông.''

Cả nhà đều im lặng, chờ đợi phần tiếp theo của câu chuyện. Ryuuhou nhìn quanh nhìn quất, thấy mọi người đều chăm chú nghe thì liền hỏi tiếp, như chính con cá trong câu chuyện đang nói: ``Nó hỏi thủ thư mượn một quyển sách, các bạn đoán là nó muốn mượn sách gì?''

Suzuno, luôn luôn lịch sự và nhẹ nhàng, liền giơ tay như đang ở trên lớp học, hỏi nhỏ: ``Là sách về biển và đại dương ạ?''

``Đúng rồi! Chính xác là sách về biển!'' Ryuuhou reo lên, nghĩ rằng câu chuyện cười của mình đã thành công.

Nhưng không ai trong số mọi người cười lên, cả bọn chỉ nhìn nhau với ánh mắt khó hiểu. Naomi thậm chí còn nghiêm túc nhíu mày suy nghĩ, sau đó ngước lên hỏi con bé với vẻ mặt thật thà: ``Chị ơi, cá mà có thể đọc được sách ạ? Sao nó vào thư viện mượn sách được?''

Kaigou bật cười lớn, vừa cười vừa nói với Ryuuhou: ``Câu hỏi của Nao mới hay hơn chuyện cười của em đấy, đúng là một câu hỏi logic mà con bé nghĩ ra.''

Ryuuhou ôm trán, rên rỉ thất vọng: ``Em biết trước mà! Biết là câu chuyện cười này sẽ bị các anh chị hiểu sai, không ai cười cả!''

Tôi liền vỗ tay hai cái lớn để cứu vãn tình thế, cười nói: ``Dù sao thì câu chuyện cũng có kết thúc tốt đẹp mà, con cá cũng tìm được đúng thư viện và mượn được đúng sách nó muốn, đâu có tệ đâu.''

Em tôi ngước lên nhìn tôi, mặt ra vẻ nghi ngờ: ``Anh đang cố cứu câu chuyện của em hay là đang làm nó trở nên tệ hơn thế nữa vậy, Onii!?''

``Anh chỉ đang cố giữ cho trò chơi của chúng ta tiếp tục thôi, chứ có ý gì đâu,'' tôi vừa cười vừa đáp lại, tay đập chan chát vào nền khăn trải bàn.

Kson cũng cười lớn, gật đầu đồng tình với lời nói của tôi: ``Tôi cho câu chuyện cười này một điểm, điểm cho sự cố gắng của Ryuuhou-chan nhé, dù không ai hiểu được điểm cười đâu!''

\ct

Ryuuhou vẫn không thể dẹp bỏ niềm tò mò về chiếc cần câu mới toanh mà mình vừa mua, đôi mắt cứ liếc đi liếc lại về phía chiếc ba lô đựng bộ dụng cụ. Nhưng Kaigou nhanh chóng nhắc nhở em một lần nữa, giọng nghiêm chỉnh: công viên này có quy định cấm mọi hoạt động câu cá để bảo vệ hệ sinh thái sông, khu vực picnic chúng ta đang ngồi cũng không ngoại lệ. 

Cpn bé thở dài nhẹ nhàng, rồi bất chợt sáng mắt ra: ``Không câu được thì lắp thử bộ dụng cụ xem sao cũng được!'' Em liền lôi chiếc hộp hướng dẫn sử dụng dày cộp ra khỏi ba lô, mở tung ra trên tấm bạt trải cỏ. Thấy có trò chơi mới, Naomi và Koneko liền nhanh nhẹn nằm sấp cạnh nhau, hai cặp mắt tròn xoe dán chặt vào những hình minh họa đầy màu sắc trong sách hướng dẫn, háo hức xem lắp ráp.

``Đầu tiên theo hướng dẫn nói là phải lắp cuộn dây câu vào cái chỗ khía này đúng không?'' Ryuuhou nói to như đang tự nhủ với mình, tay cầm chiếc cuộn dây nhựa đen thẳng thừng, định gắn vào thân cần.

Kaigou ngồi kế bên, tay vừa gấp gọn chiếc khăn ăn vừa liếc qua hình minh họa tiếp theo trong sách, lập tức lắc đầu chỉ vào bức tranh: ``Em dòm kỹ đi, cuộn dây của em đang lắp ngược rồi kìa. Mặt có rãnh luồn dây phải quay ra ngoài chứ.''

Ryuuhou vội vàng xoay ngược chiếc cuộn dây lại, đối chiếu với hình trong sách, rồi bật cười ngượng: ``Ối đúng là ngược thật, may mà chị chỉ kịp.''

``Chị đã biết em sẽ làm sai từ lúc cầm cuộn dây lên rồi mà,'' cô chị Hai trêu chọc.

``Em đang tự kiểm tra xem mình có phát hiện ra lỗi không đâu, không phải lúng túng gì đâu!'' Ryuuhou nhanh chóng biện minh, tay vẫn chỉnh lại vị trí cuộn dây cho chắc chắn.

Tôi cũng bước lại ngồi xuống cạnh ba đứa trẻ, nhìn vào bộ dụng cụ đang được dần lắp ráp. Nhìn em ấy đang cố gắng luồn sợi dây qua các rãnh nhỏ, tôi liền góp ý: ``Để cắt dây cho đúng kích thước thì em cần thêm một cái kìm nhỏ ở nhà đấy, tay thường không cắt nổi sợi cước thép này đâu.''

Ryuuhou ngước đầu lên khỏi chiếc cần câu, mặt đầy ngạc nhiên: ``Anh cũng biết cách lắp những thứ này à?''

``Anh hồi trước có từng sửa vài món đồ hỏng trong nhà, nên cũng rành chút ít về dụng cụ,'' tôi nhún vai đáp.

Con bé lập tức nở nụ cười rạng rỡ, như thể vừa nghĩ ra một kế hoạch tuyệt vời: ``Thế thì lần sau chúng ta đi câu ngoài thành phố, anh phải đi cùng bọn em nhé! Anh mà có đi thì em mới yên tâm.''

Kaigou lập tức khoanh tay trước ngực, ra điều kiện ngay: ``Đi được cũng được, nhưng trước hết em phải học cách buộc các nút dây câu cho đúng, học thuộc lòng các kỹ thuật cơ bản đã.''

``Được thôi, em hứa sẽ học thật siêng!'' Ryuuhou nhanh chóng đồng ý, tay lại tiếp tục luồn dây qua các khớp của cần câu.

Naomi lúc này nhặt lên một đoạn dây cước thừa, đưa lên trước mặt ngắm nghía. Ánh nắng chiều chiếu vào sợi dây, làm nó phản chiếu những vệt sáng bạc lung linh. Em nhỏ giọng kêu lên: ``Sợi dây này mảnh như tơ vải vậy, nhìn yếu ớt thật.''

Koneko cũng chạm nhẹ ngón tay vào cuộn dây còn lại trên tấm bạt, sờ thử độ dai của nó, rồi hỏi nhỏ Ryuuhou với vẻ tò mò: ``Nó có dễ bị đứt khi móc phải vật nặng hay cá to không ạ?''

``Kéo đúng cách và buộc nút đủ chặt thì không bao giờ đứt đâu,'' con bé giải thích thật kỹ, giọng nói chậm rãi, rõ ràng, hoàn toàn không giễu cợt hay tỏ vẻ sốt ruột trước những câu hỏi đơn giản của hai đứa nhỏ. Naomi chăm chú nghe, gật đầu liên tục để ghi nhớ, và Koneko cũng vậy, hai cặp mắt tròn xoe nhìn cô chị như những học trò nhỏ đang nghe cô giáo giảng bài.

Sau vài phút bận rộn, chiếc cần câu cuối cùng cũng được lắp xong xuôi. Ryuuhou cầm thử lên tay, vung nhẹ một cái để kiểm tra độ cân bằng, rồi lại cẩn thận tháo dỡ ra, cuộn lại từng phần để cất vào hộp. ``Hôm nay chỉ học lắp ráp thôi là đủ rồi, ngày mai tiếp tục học các bước tiếp theo,'' em nói với mọi người.

Kaigou nhìn em ấy cất giữ bộ dụng cụ một cách cẩn thận, gật đầu đầy hài lòng: ``Biết dừng đúng lúc, không tham lam làm quá sức mình, đó cũng là một phần quan trọng của việc học bất cứ kỹ năng mới nào đấy.''

Ryuuhou ngước lên, miệng nở nụ cười tinh nghịch: ``Chị đang khen em đấy à?''

``Chị chỉ đang ghi nhận sự tiến bộ của em thôi, không có khen gì đâu,'' Kaigou vờn mặt lạnh, nhưng khóe miệng vẫn nhô lên một chút.

``Dù sao em cũng sẽ coi đây là lời khen, và sẽ cố gắng hơn nữa!'' Ryuuhou cười lớn, đóng chặt nắp hộp đựng cần câu lại.

\ct

Thời gian trôi qua thật nhanh hơn bất kỳ ai trong chúng tôi dự đoán, những phút giây vui vẻ cứ thế trôi đi lặng lẽ đến nỗi chẳng ai nhận ra bầu trời đã dần chuyển sắc. Ánh nắng buổi trưa vốn còn gắt gỏng giờ cũng mềm lại, chếch dần về phía chân trời phía tây, vắt những vệt vàng ấm áp lên những tán lá cây cổ thụ trong công viên. Giữa khoảng bãi cỏ xanh mướt, Naomi và Koneko đã rủ nhau chơi chuyền bóng với quả bóng nhỏ màu cam mà chúng tôi mới mua lúc nãy.

Lần đầu tiên cầm bóng, con bé còn ngần ngại ném quá mạnh, chỉ đưa nhẹ quả bóng về phía Koneko với một lực rất vừa phải. Nhưng em ấy vẫn chưa quen với tốc độ của quả bóng, hai tay dang rộng mà bắt hụt lần đầu, bóng lăn bật trên mặt cỏ mịn chỉ vài bước rồi dừng lại. Naomi lập tức chạy nhảy về phía quả bóng, cúi xuống nhặt lên với nụ cười tươi rói.

``Xin lỗi nhé, tớ ném hơi nhẹ quá,'' em gọi với về phía cô bạn, giọng đầy vẻ ngại ngần.

Koneko lắc đầu, mái tóc ngắn bay nhẹ theo chuyển động: ``Không sao đâu, tớ cũng chưa quen với việc bắt bóng thôi, lần sau chắc chắn sẽ được.''

Đến lần thứ hai, Naomi điều chỉnh lực ném vừa đủ, quả bóng bay thẳng về phía cô bạn tai mèo với tốc độ dễ chịu hơn. Lần này Koneko đã chuẩn bị sẵn, hai tay chụp lấy quả bóng một cách chắc chắn, ôm nó vào lòng để không rơi ra. 

Naomi thấy vậy liền vỗ tay reo lên mừng rỡ, tiếng cười trong trẻo vang vọng khắp khoảng không gian quanh đó: ``Cậu bắt được rồi! Tuyệt quá Koneko-chan!''

Koneko cũng nở ra nụ cười rạng rỡ, đôi mắt sáng bừng lên vì niềm vui nhỏ bé ấy. Em ném quả bóng trở lại về phía Naomi, nhưng có lẽ vì còn hơi hồi hộp nên quả bóng bay hơi lệch sang một bên. Naomi nhanh nhẹn chạy sang hướng đó, nhảy lên bắt lấy quả bóng chỉ cách mép bãi cỏ vài bước, gần kề với lối đi nhỏ xung quanh cây cổ thụ. 

Từ đó về sau, hai đứa cứ thế chơi đi chơi lại, không ai bận tâm đếm điểm hay đặt ra những quy tắc phức tạp gì cả. Chỉ có quả bóng nhỏ màu cam cứ thế bay qua bay lại giữa hai cô bé, mang theo niềm vui đơn giản mà trong lành của tuổi thơ.

Tôi ngồi dưới bóng cây cùng Delutaya, ánh mắt dõi theo hai đứa trẻ đang chơi đùa tự do. Cô nàng khẽ mỉm cười, giọng nói nhỏ nhẹ như sợ làm vỡ cảnh tượng êm dịu trước mắt: ``Koneko-chan hôm nay cười nhiều hơn hẳn so với lúc sáng mới đến nhỉ? Em thấy em ấy mở lòng hơn nhiều rồi.''

Tôi gật đầu đồng tình, mắt vẫn không rời khỏi hai bóng nhỏ đang nhảy nhót trên bãi cỏ: ``Naomi cũng vậy. Từ lúc đến đây, em ấy chưa bao giờ ngừng cười cả. Hai đứa hợp nhau thật sự.''

Delutaya tiếp tục: ``Em nghĩ từ nay về sau, hai đứa sẽ còn nhiều cơ hội để đi chơi cùng nhau nhiều lần nữa. Chúng ta có thể sắp xếp những chuyến đi như thế này thường xuyên hơn.''

``Anh cũng mong là thế,'' tôi đáp lại, lòng tràn đầy sự ấm áp khi thấy hai đứa trẻ hòa hợp với nhau.

Lúc này Kson từ phía bên kia bò lại, tay cầm hai chai nước lạnh lấy từ túi giữ nhiệt, đưa cho tôi một chai. Cô ngồi xuống cạnh tôi, lau nhẹ mồ hôi trên trán rồi nói: ``Đi chơi thế này này mới đúng là nghỉ ngơi chứ! Không cần phải lên kế hoạch chi tiết tới từng phút, không ai phải thúc giục ai làm gì, cứ tự nhiên mà tận hưởng khoảnh khắc thôi!''

Tôi cười, nhận lấy chai nước từ tay cô: ``Anh có lên kế hoạch mà, từ ngày hôm trước đã lên lịch trình đầy đủ rồi. Chỉ là chờ thời điểm thích hợp thôi.''

Nàng tóc tím cười khúc khích, lắc đầu: ``Em biết chứ, Kuon-paisen lúc nào cũng chu đáo như vậy. Nhưng điều hay là anh vẫn để mọi người tự do chọn xem mình muốn làm gì, không bắt ai phải theo đúng kế hoạch một cách cứng nhắc.''

Tôi nhìn cô, rồi liếc qua những người bạn đang tụ tập quanh bãi cỏ, giọng nói đầy nhẹ nhàng: ``Hôm nay anh không phải quản lý ai cả. Chỉ là một người bạn cùng đi chơi thôi.''

Kson bật cười lớn, tiếng cười vang lên khiến mấy người khác cũng phải quay sang nhìn. Cô chỉ vào tôi, giọng đầy trêu chọc: ``Anh cứ tự thuyết phục mình như thế đi, ai mà chẳng biết anh là người lo lắng cho mọi thứ nhất trong nhóm đâu!''

Đúng lúc đó, từ phía sau tôi vang lên tiếng gọi vọng tới của Kaigou, cô vừa từ bên kia bãi cỏ đi về, mắt thấy chai nước trong tay tôi liền la lên: ``Aniki! Đừng có uống nhầm nước của em! Cái đó là chai của em đấy!''

Tôi ngạc nhiên nhìn chai nước trong tay, lật nắp lên xem thì thấy đúng là có dán nhãn tên Kaigou bằng bút ghi nhớ màu hồng mà em ấy hay dùng. Tôi cười ngượng, đưa chai nước trả lại cho cô: ``Anh tưởng Kson đưa chai của anh, ai ngờ lấy nhầm. Xin lỗi nhé.''

Kson liền lục lại trong túi giữ nhiệt, lấy ra một chai nước khác có dán tên tôi rồi đưa qua: ``Ối xin lỗi anh, em lấy nhầm. Đây mới là chai của anh đây, đúng nhãn tên rồi.''

Tôi nhận lấy chai nước, mở nắp uống một ngụm mát lạnh: ``Cảm ơn em nhé.''

Kaigou nhanh nhẹn giật lại chai nước của mình, ôm vào lòng như sợ tôi lại lấy nhầm lần nữa. Cô nhấn mạnh, giọng đầy vẻ cảnh báo: ``Lần sau nhìn nhãn tên trước khi uống nhé, đừng có mà uống nhầm đồ của người khác nữa. Em đã ghi tên rõ ràng thế mà cũng lấy nhầm!''

\ct

Trước khi bước chân rời khỏi công viên, Suzuno đã đứng dậy làm nhiệm vụ kiểm tra lại toàn bộ đồ đạc, đảm bảo không có món nào bị quên sót lại trên bãi cỏ. Túi rác chứa đầy vỏ hộp và giấy gói đã được buộc chặt miệng, những hộp đựng thức ăn thừa được xếp gọn gàng vào túi giữ nhiệt để giữ được độ tươi ngon. Tấm bạt trải cỏ lớn được phủi sạch hết những cọng cỏ và lá cây rơi vãi, sau đó được gấp gọn vào trong một chiếc túi vải. Naomi chạy nhảy khắp nơi để nhặt lại quả bóng màu cam, còn Koneko thì nhẹ nhàng gấp những chiếc khăn giấy còn sạch sẽ để có thể tái sử dụng. 

Ryuuhou nhanh nhẹn đeo chiếc ba lô to đựng bộ cần câu mới mua vào lưng, còn Kaigou thì xách chiếc hộp chứa các đồ gia dụng đã dùng trong buổi dã ngoại. Tôi cầm chiếc túi đựng những thực phẩm còn thừa không thể bỏ vào túi giữ nhiệt, còn Kson và Delutaya thì cùng nhau chia sẻ chiếc túi nhỏ đựng những chiếc thẻ nhớ vừa mua và vài món phụ kiện nhỏ lẻ khác.

``Chúng ta còn cần phải ghé qua cửa hàng tiện lợi nào không ạ?'' nàng Tam giác vừa xách túi vừa hỏi với giọng thắc mắc.

Suzuno mở chiếc danh sách đã ghi từ trước, lướt qua từng mục một cách cẩn thận. ``Không cần đâu ạ. Chúng ta đã mua đủ tất cả mọi thứ cần thiết rồi.''

Nghe vậy, Ryuuhou đột nhiên nhìn lại tờ danh sách trên tay, rồi ngẩng lên nói với vẻ nghĩ ra gì đó quan trọng: ``Em lại còn thiếu một món nữa.''

Kaigou lập tức quay sang nhìn em gái với ánh mắt cảnh giác. ``Món gì nữa?''

``Mồi câu ạ.''

``Mơ đi,'' cô em song sinh lắc đầu ngay lập tức.

``Em chỉ hỏi thôi mà chị!'' Ryuuhou cười trừ, vẫy tay để làm dịu bầu không khí.

Naomi nhìn vào chiếc ba lô đựng bộ cần câu trên lưng Ryuuhou, rồi ngẩng lên hỏi với giọng tò mò: ``Mình có thể đi tìm mồi ở gần bờ sông không ạ? Chắc có nhiều côn trùng ở đó phải không?''

``Không được đâu em, công viên có quy định không được tự ý bắt côn trùng hay động vật nhỏ trong công viên đâu,'' cô chị Hai nhẹ nhàng giải thích cho em ấy hiểu.

Ryuuhou nghe vậy cũng chỉ chẹp miệng chấp nhận, rồi nói với vẻ đã chuẩn bị sẵn kế hoạch: ``Vậy lần sau em sẽ mua mồi câu ở cửa hàng chuyên dụng về, không làm phiền gì đến côn trùng ở đây đâu.''

``Lần sau nếu muốn mua gì cũng phải hỏi trước chị và mọi người đã, đừng tự ý mua lung tung đấy!'' Kaigou nhấn mạnh một lần nữa.

``Em biết rồi, lần sau em sẽ hỏi trước mà!'' con bé đáp nhanh chóng, vừa đi vừa tung tăng theo sau mọi người.

Chúng tôi cùng nhau bước đi về phía trạm xe buýt, bước chân nhẹ nhàng trên con đường mòn quanh co trong công viên. Naomi và Koneko đi sát cạnh nhau, không ngừng chuyền quả bóng màu cam qua lại bằng đôi tay nhỏ bé của mình. Mỗi lần quả bóng bay lệch hướng, cả hai đứa lại cùng nhau cười lớn rồi chạy đến nhặt lên, tiếng cười trong trẻo vang vọng khắp không gian. Tôi nhìn theo bóng dáng hai đứa trẻ, lòng cảm thấy nhẹ nhàng và bình yên đi rất nhiều, mọi mệt mỏi trong ngày dường như tan biến hết.

Mikeneko hôm nay đã không đi cùng buổi dã ngoại, nhưng em ấy chắc chắn biết rằng chúng tôi sẽ sớm trở về nhà. Có lẽ với một ngày đầy đủ niềm vui và bình yên như hôm nay, chỉ cần vậy là đủ rồi.

\ct

Chuyến xe buýt trên đường về nhà đông đúc hơn hẳn so với lúc chúng tôi đi vào buổi sáng. Chúng tôi chọn đứng ở phía cuối xe, tất cả các túi đồ đã được xếp gọn gàng dưới chân để không cản trở lối đi. Ryuuhou luôn giữ chiếc ba lô đựng bộ cần câu sát trước ngực, sợ rằng nó sẽ bị va đập vào người khác trên xe đông đúc, còn Kaigou thì đứng cạnh cửa xe để tiện cho việc xuống xe khi đến trạm. 

Suzuno may mắn có được một ghế ngồi cạnh cửa sổ, và em ấy mời Naomi cùng Koneko ngồi cạnh mình để hai đứa trẻ không phải đứng mệt. Kson và Delutaya đứng ở phía sau chúng tôi, không ngừng nói nhỏ với nhau về chiếc mô hình nhân vật mà họ vừa thấy ở cửa hàng, đôi mắt sáng rực lên khi nhắc đến món đồ mình thích. Tôi thì giữ chặt chiếc hộp thủy tinh nhỏ bằng hai tay, đó là món đồ nhỏ tôi mua ở cửa hàng lưu niệm của công viên để để ở phòng khách.

Mỗi khi tài xế phanh xe gấp, tất cả mọi người trên xe đều khẽ nghiêng người theo quán tính, có người còn phải bám chặt vào tay vịn để tránh bị ngã. Có một lần xe phanh đột ngột, Naomi suýt nữa làm rơi quả bóng màu cam đang giữ trong tay, may mắn là Koneko đã nhanh tay đỡ lấy nó trước khi quả bóng lăn ra lối đi giữa xe. Cô em út ngước lên nhìn bạn mình với ánh mắt biết ơn, còn bé Mèo con chỉ nhẹ nhàng gật đầu, rồi giữ quả bóng ấy trong lòng cho đến lúc chúng tôi xuống xe. Khi xe dừng lại ở trạm đối diện nhà chúng tôi, mặt trời đã thấp dần xuống phía chân trời, nhuốm màu vàng ấm áp lên toàn bộ khu phố.

Vừa bước chân xuống khỏi xe buýt, Ryuuhou đã liền quay sang hỏi tôi với giọng thắc mắc: ``Bây giờ là mấy giờ rồi anh?''

``Bốn giờ mười phút rồi,'' tôi nhìn vào đồng hồ trên tay rồi trả lời.

``Vậy thì chúng ta còn kịp về nhà trước khi trời tối, may quá,'' con bé thở phào nhẹ nhõm.

Kaigou nhìn vào đống túi đồ mà chúng tôi đang xách, rồi lắc đầu nói với em gái: ``Chúng ta có đi xa đến đâu đâu mà phải lo về trời tối? Ngay trước mặt nhà mà.''

``Nhưng em mệt lắm rồi,'' Ryuuhou rên rỉ, dường như mọi năng lượng đã cạn kiệt sau một ngày vui chơi.

``Em mệt là vì suốt ngày cứ đứng ngắm nghía bộ cần câu mới mua thôi, đâu có làm gì nặng nhọc đâu,'' cô trêu chọc em gái, khiến mọi người xung quanh đều bật cười.

``Ngắm đồ mới mua cũng tốn sức mà, phải nghĩ xem ngày mai sẽ dùng nó như thế nào nữa chứ!'' Ryuuhou nhanh chóng bao biện, đáp lai chỉ là cái lườm nguýt từ Kaigou.

Tôi bước lên mở cửa nhà, nhường tất cả mọi người vào trước. Koneko đột nhiên quay lại nhìn tôi, giọng nói nhỏ nhẹ nhưng đầy chân thành: ``Cảm ơn Onii-chan vì buổi đi chơi hôm nay ạ.''

``Có gì đâu mà cảm ơn. Hôm nay em có vui không?'' tôi mỉm cười hỏi em.

Koneko gật đầu lia lịa, đôi mắt sáng rực lên vì niềm vui. Naomi đứng cạnh em, vẫn ôm chặt quả bóng màu cam trong người, rồi cũng nói thêm với giọng háo hức: ``Lần sau chúng ta lại đi chơi như thế này nữa nhé, được không ạ?''

Koneko nhìn vào gương mặt háo hức của Naomi, rồi mỉm cười và đáp lại một từ ngắn gọn nhưng đầy nhiệt tình: ``Ừ. Nếu mà anh có dịp.''

\ct

Khi màn đêm dần buông xuống căn nhà chung của chúng tôi, ánh đèn vàng ấm áp của căn bếp ở tầng năm đã bật sáng từ rất sớm, lan tỏa sự ấm cúng khắp không gian sau một ngày dài vui chơi ngoài trời. Tôi nhanh chóng thay bộ quần áo dính chút bụi cỏ từ công viên, rửa sạch tay chân rồi bắt tay vào chuẩn bị bữa tối cho cả nhà. 

Suzuno, với sự chu đáo thường ngày, đã rửa sạch từng bó rau xanh và nấu sẵn nồi cơm thơm phức từ trước, chỉ chờ các món ăn khác hoàn thành là có thể dọn lên bàn. Kaigou ngồi cạnh bếp, giúp tôi thái những củ hành tây thơm lừng thành những lát mỏng đều đặn, trong khi Ryuuhou thì đứng dựa vào góc bàn, mắt vẫn không rời khỏi cuốn sách hướng dẫn lắp ráp bộ cần câu mới mua của mình, nhẩm đi nhẩm lại từng bước như sợ quên. 

Phòng khách cũng không kém phần náo nhiệt: Naomi và Koneko ngồi xếp bằng trên thảm, tập trung tô màu bức tranh vẽ cảnh công viên mà chúng tôi vừa ở, những nét bút màu rực rỡ dần che đi phần giấy trắng. Còn Kson thì ngồi trên ghế sofa, say sưa kể cho Delutaya nghe về chiếc mô hình nhân vật có thể đổi nhiều tư thế mà cô thấy ở cửa hàng lưu niệm, giọng đầy sự hấp dẫn.

``Chị vẫn nghĩ mình nên quay lại mua nó thôi. Nó quá hợp với bộ sưu tập của chị,'' nàng Tổng trưởng lặp lại lần nữa, mắt vẫn lấp lánh sự mong muốn.

``Chị đã nói đúng vậy bốn lần từ lúc về nhà rồi đấy,'' nàng Tam giác đáp lại với giọng cười khúc khích, biết cô bạn mình vẫn không dứt được niềm đam mê với món đồ đó.

``Vì nó vẫn là một ý đúng mà! Ai mà không muốn có thêm một chiếc mô hình đẹp trong phòng chứ?''

Tôi vừa khuấy đều nồi canh đang sôi lăn tăn trên bếp, vừa vô tình liếc nhìn đồng hồ treo tường. Mikeneko, người đã ở nhà suốt ngày hôm nay, vẫn chưa bước xuống tầng năm. Tôi liền mở tủ bát, lấy thêm một chiếc bát gốm sạch để dành phần cho em ấy. 

Kaigou, vừa tiếp tục thái hành vừa nhìn thấy hành động của tôi, liền hỏi nhỏ: ``Anh đang để phần cho Mike-chan à?''

``Ừ, chắc em ấy cũng sắp xuống thôi.''

``Có cần em gọi điện hay lên phòng gọi em ấy xuống ăn cơm không? Để tránh phần ăn bị nguội,'' Kaigou đề xuất, đặt con dao thái rau xuống.

``Không cần đâu. Chờ cơm hoàn toàn dọn xong, anh sẽ gọi. Không cần phải vội,'' tôi lắc đầu từ chối, biết rằng Mikeneko thường thích có không gian riêng của mình trước khi tham gia vào bữa ăn chung.

Kaigou chỉ gật đầu hiểu ý, không hỏi thêm gì nữa; chúng tôi đều quen với thói quen của em ấy rồi. Dần dần, mùi thơm của các món ăn vừa nấu xong lan tỏa khắp toàn bộ căn phòng khách, lôi kéo sự chú ý của mọi người. Suzuno là người đầu tiên mang tô canh lớn nóng hổi đặt giữa bàn ăn. Naomi và Koneko nghe vậy liền nhanh chóng cất những cây bút màu vào hộp, đứng dậy cùng nhau xếp đũa thìa cho tất cả mọi người, sự háo hức hiện rõ trên mặt. Ryuuhou cũng đóng cuốn sách hướng dẫn lại, cất vào hộp đựng cần câu của mình rồi bước đến bàn.

Delutaya đem đĩa rau luộc thanh mát đặt cạnh tô canh, còn Kson thì mở rộng cửa ban công để mùi thức ăn từ bếp không bị đọng lại trong nhà. Tôi bắt đầu chia từng phần thức ăn vào bát của mỗi người, và một phần được đặt riêng trong một chiếc bát nhỏ, để dành cho Mikeneko. Sau khi dọn xong, tôi không khỏi ngước nhìn lên cầu thang dẫn lên tầng tám, nơi phòng của em ấy nằm.

Mọi người bắt đầu ngồi xuống ăn trước, tiếng cười nói rộn ràng vang lên. Naomi là người hào hứng nhất, liên tục kể lại cho mọi người nghe khoảnh khắc Koneko bắt thành công quả bóng màu cam ở công viên, khiến bé Mèo con phải cúi đầu, hai má ửng đỏ vì ngại ngùng trước sự khen ngợi.

``Cậu ấy bắt bóng giỏi lắm! Lần đầu tiên tập mà làm được thế là tuyệt vời rồi,'' Naomi nói liên tục, tay chỉ vào bạn mình để nhấn mạnh.

Ryuuhou cũng gật gù đồng tình, miệng vẫn nhai thức ăn: ``Ừ nhỉ, lần sau chúng ta có đi chơi đâu đó, chị em mình lập đội chơi chuyền bóng nhé. Chắc chắn sẽ thắng hết mọi người.''

Koneko ngước nhìn Ryuuhou, hơi bất ngờ trước lời mời đó, mắt tròn xoe một chút.

``Nhưng phải là nếu cậu ấy muốn thôi nhé,'' Naomi nhanh chóng bổ sung, sợ bạn mình cảm thấy bị ép buộc.

``Dĩ nhiên rồi, nếu Koneko-chan muốn thì cứ nhắn tin gọi em là được. Em luôn sẵn sàng lập đội,'' Ryuuhou mỉm cười đảm bảo, làm cho gương mặt của Koneko cũng nở ra nụ cười nhẹ nhõm.

Đầu bữa ăn, Suzuno còn chủ động gắp thêm một đũa rau xanh vào bát của Kaigou, khiến cô phải nhìn lượng rau đã chất đầy trong bát mình rồi lại nhìn cô em gái của mình.

``Suzu, chị ăn được rồi đó. Đủ rau lắm rồi,'' Kaigou than thở, biết rằng Suzuno luôn quan tâm đến mình dù đã bước sang cuộc sống sung túc hơn.

``Em biết ạ, Onee-sama. Nhưng đây mới đúng là phần rau chị cần ăn cho đủ dinh dưỡng hôm nay,'' Suzuno đáp lại một cách rất nghiêm túc, không nhượng bộ.

``Em có thể nói lời đó nhẹ nhàng hơn một chút không? Giống như mời chị ăn thay vì ra lệnh thế này?''

``Onee-sama cần ăn rau ạ,'' Suzuno lặp lại đúng câu đó, chỉ thay đổi cách xưng hô, nhưng giọng vẫn vô cùng chắc chắn.

``Nghe vẫn y hệt một mệnh lệnh từ thầy hiệu trưởng thế nào ấy,'' Kaigou thở dài giả vờ bất lực, khiến mọi người xung quanh đều bật cười.

Trong khi đó, Kson vừa ăn vừa say sưa kể lại chuyện hai người bị lạc trong cửa hàng điện tử lúc chiều, còn Delutaya thì liên tục đính chính lại câu chuyện của cô bạn, rằng họ chỉ bị lạc khỏi quầy thẻ nhớ thôi chứ không hề lạc ra khỏi cả cửa hàng như Kson nói.

``Hai người đi lạc đến mức đã đi qua hết ba tầng của cửa hàng đó rồi mà còn chối,'' tôi nhắc vào, nhớ lại lúc cả nhóm phải đi tìm hai người trong cửa hàng.

``Đó là vì cửa hàng có có ba tầng mà, đi qua hết là chuyện bình thường thôi,'' Kson biện minh.

``Cửa hàng đó có năm tầng cơ, hai người mới chỉ đi được hơn nửa thôi,'' tôi phá vỡ lời biện minh của cô, khiến cả bàn lại cười lớn.

Ryuuhou bật cười thành tiếng, suýt nữa làm sôi cơm trong miệng. Cả buổi ăn trôi qua rất vui vẻ, chỉ có chiếc ghế đối diện tôi - chỗ ngồi thường ngày của Mikeneko - vẫn còn trống không. Tôi đợi cho mọi người ăn được một nửa bữa, không khí đã thật sự thoải mái, thì mới đứng dậy khỏi ghế.

``Anh lên phòng gọi Mike-chan một tiếng, bảo em ấy xuống ăn thôi.''

Kaigou ngồi cạnh tôi liền gật đầu, nhắc nhở nhỏ: ``Cứ nói với em ấy là còn phần ăn, cơm canh vẫn còn nóng. Không cần phải giải thích hay hỏi thêm gì đâu, để em ấy tự nhiên.''

``Anh biết rồi, yên tâm.''

\ct

Tôi bước lên cầu thang với những bước chân khẽ nhũn, cố gắng giữ tiếng dép lướt trên bậc gỗ thật nhỏ, sợ rằng bất kỳ tiếng động nào cũng sẽ làm gián đoạn không khí ấm áp đang vây lấy cả căn nhà. Vừa đến chiếu nghỉ giữa tầng năm và tầng sáu, tai tôi bỗng bắt được tiếng cửa phòng ở tầng trên kêu rất nhẹ - một tiếng mở cửa vừa đủ để một người lách ra, không đủ để làm ai khác chú ý. 

Ngước mắt lên, tôi thấy Mikeneko đứng ở đầu cầu thang tầng tám, thân người vẫn còn gần như nửa ẩn trong khoảng tối của căn phòng mình. Một bàn tay em vẫn nắm chặt mép cửa, như thể vẫn chưa hoàn toàn sẵn sàng bước ra khỏi không gian riêng ấy. Ánh mắt em không nhìn tôi, mà hướng về phía phòng ăn phía dưới, nơi tiếng cười nói rộn ràng của mọi người vẫn đang vọng lên, len lỏi qua từng tầng cầu thang.

Tôi dừng lại ngay, không bước thêm một bậc nào, giữ khoảng cách chỉ vài bậc thang giữa chúng tôi, vừa đủ để em cảm thấy thoải mái, không bị ép buộc phải tiếp cận.

``Bữa tối chưa xong đâu,'' tôi nhỏ giọng nói, sửa lại câu nói đầu tiên thành một lời thông báo nhẹ nhàng, không vội vàng. ``Mọi người vẫn đang ăn.''

Mikeneko cuối cùng cũng chuyển mắt xuống, nhìn vào chiếc bát gốm sạch tôi đang cầm bằng cả hai tay, hơi nước còn bay nghi ngút từ miệng bát. ``Anh mang gì lên đấy?'' em hỏi, giọng cũng không lớn hơn tiếng thì thầm.

``Phần của em. Anh để riêng cho em từ lúc dọn cơm.''

``Em đâu có bảo anh phải để phần đâu,'' em đáp lại, có chút ngạc nhiên lẫn lưỡng lự.

``Anh biết. Nhưng vẫn muốn để thôi.''

Em im lặng một lúc dài, chỉ có tiếng cười của Naomi từ dưới nhà vọng lên phá vỡ khoảng trống giữa chúng tôi. Tôi không vội, chỉ đứng yên đó rồi từ từ đặt chiếc bát lên bậc thang rộng bên cạnh tôi, không tiến thêm một bước nào về phía em. Tôi muốn để em có toàn quyền chọn lựa, không có bất kỳ áp lực nào.

``Nếu em cảm thấy ổn khi ăn dưới nhà thì mọi người vẫn còn ngồi đó hết. Còn nếu em muốn ăn ở đây một mình, anh để phần này lại cũng được. Anh không nói gì thêm với ai đâu.''

Mikeneko nhìn chiếc bát đặt trên bậc thang, rồi lại ngước mắt nhìn tôi, ánh mắt trong phòng tối bỗng có chút ấm lại. ``Hôm nay mọi người đi chơi ở công viên vui không?'' em hỏi, câu hỏi như thể đã chờ để nói từ lâu.

``Vui lắm. Naomi với Koneko-chan suốt ngày chơi chuyền bóng, không bao giờ chán. Ryuuhou thì mua cả một bộ cần câu mới, đang mơ đi câu ngoài thành phố.''

``Em ấy biết cách câu cá đâu?'' Mikeneko khẽ nhíu mày, có chút hoài nghi.

``Chưa biết được. Kaigou vừa mới hứa sẽ chỉ em ấy cách buộc nút dây trước, chưa hứa gì về việc đi câu cả.''

Cô Mèo hòng bật cười khúc khích, tiếng cười nhỏ nhẹ như gió thổi qua lá cây. ``Nghe đúng kiểu Ryuuhou-chan rồi, chắc là sắp mua hết cửa hàng dụng cụ câu cá về nhà mất.''

``Anh cũng vừa mới ngăn cản được em ấy mua mồi câu hôm nay, may mà kịp.''

``Vậy là em cứu được ví tiền của anh rồi nhỉ,'' em trêu, khóe miệng nhô lên một nụ cười nhỏ.

``Anh sẽ ghi nhận công lao to lớn đó vào sổ công của nhà.''

Vừa lúc đó, tiếng la hét của Ryuuhou từ dưới nhà vọng lên, làm cả cầu thang đều rung nhẹ: ``\textbf{Ai ăn chậm là hết phần kem tráng miệng nhé! Em ăn hết cái hũ matcha đây!}'' Mikeneko nghiêng đầu, như thể không thể tin vào những gì mình vừa nghe.

``Nó đang dọa ai thế này?''

``Chắc chắn là Kson rồi, có lẽ nó lại đi rửa tay lâu quá nên sợ mất phần.''

Mikeneko thở dài một tiếng, rồi cuối cùng cũng buông mép cửa ra, bước nhẹ xuống vài bậc thang. Nhưng em lại dừng lại giữa chừng, như thể còn cần thêm chút thời gian để chuẩn bị hòa vào không khí náo nhiệt dưới kia. Tôi không giục em, chỉ đứng yên đó chờ. Sau vài giây im lặng, em lại tiếp tục bước xuống, chân dép khẽ chạm vào từng bậc thang. 

Từ phòng ăn, Koneko là người đầu tiên thấy chị mình xuất hiện ở đầu cầu thang. Em bé không nói gì, chỉ lẳng lặng dịch chiếc ghế trống cạnh mình ra xa một chút, để người chị có đủ chỗ ngồi thoải mái. 

Mikeneko bước đến, ngồi xuống cạnh em, hai vai em dần dần thả lỏng ra. Không ai hỏi em tại sao hôm nay lại không đi cùng buổi dã ngoại. Không ai ép em phải kể lại mình đã làm gì ở nhà suốt ngày. Chỉ có Suzuno, với sự chu đáo quen thuộc, nhẹ nhàng đặt chiếc bát cơm nóng hổi ngay trước mặt em, không một lời phàn nàn hay thắc mắc.

``Mời chị dùng bữa ạ, Mikeneko-san'' Suzuno nhỏ giọng nói.

Nàng Mèo hồng cầm đôi đũa tre trên bàn lên, ngón tay khẽ chạm vào mặt gỗ mịn. ``Cảm ơn em,'' em đáp lại, giọng đã hoàn toàn bình thường.

Tôi cũng quay về chỗ ngồi cũ của mình, lắng nghe tiếng cười nói lại trở nên rộn ràng như trước. Chuyến đi chơi của cả nhà hôm nay đã kết thúc từ lâu, mọi mệt mỏi đều tan biến. Nhưng trong bữa tối ấm áp này, vẫn luôn có một chỗ trống dành cho những người đến muộn; những người cần thêm chút thời gian để bước ra khỏi không gian riêng của mình, để hòa vào cái ấm áp của gia đình.

%==========================================TRUYỆN PHỤ=========================================

\omk

\pov{-}

Sáng hôm sau, ánh nắng vàng ươm vừa len lỏi qua cửa sổ phòng khách, Ryuuhou đã hớn hở đem chiếc hộp đựng bộ cần câu mới mua đặt lên mặt bàn gỗ. Cô em song sinh Kaigou đứng ngay phía đối diện, hai tay khoanh chặt trước ngực với vẻ mặt nghiêm chỉnh như đang chuẩn bị kiểm tra một bài học quan trọng.

``Hôm nay chị sẽ giữ lời, dạy em từng bước cách buộc các nút dây câu cơ bản.''

Ryuuhou mừng rỡ reo lên: ``Chị thật sự nhớ lời hứa rồi đấy! Em tưởng chị quên mất rồi chứ!''

Kaigou lập tức nhấn mạnh để tránh em gái hiểu lầm: ``Nhưng nhớ nhé, chị chỉ nói sẽ chỉ em cách buộc dây thôi. Chưa hứa là sẽ đi cùng em ra bờ sông câu cá đâu.''

Cô em gái vội gật đầu lia lịa, mắt vẫn không rời khỏi cuốn hướng dẫn trong hộp: ``Em biết mà, em chỉ muốn học kỹ thuật trước thôi.''

``Còn một điều nữa,'' Kaigou tiếp tục liệt kê các quy tắc, ``tất cả đoạn dây thừa sau khi buộc xong, em phải tự dọn dẹp sạch sẽ, không được để vương vãi khắp nhà.''

``Em biết rồi, sẽ dọn ngay mà.''

``Không được để Naomi nhỏ đụng vào lưỡi câu, quá sắc bén dễ bị thương.''

``Em sẽ cất kỹ, không để em ấy chạm vào đâu.''

``Và tuyệt đối không được vung cần câu lung tung trong nhà, làm vỡ đồ đạc thì em chịu trách nhiệm.''

Lúc này Ryuuhou chỉ còn cách ngửa đầu nhìn trần nhà than thở: ``Em đâu có định làm những điều ngốc nghếch đó đâu, Kaigou-nee\~{}''

Hài lòng với những lời hứa của em gái, Kaigou rút từ trong túi ra một cuộn dây cước mới cứng, đưa cho con bé: ``Tốt lắm. Nếu đã hiểu hết thì bắt đầu thôi.''

Vừa lúc đó, Kson đi lè nhè ngang qua phòng khách, mắt vẫn nhắm hờ vì mới ngủ dậy, tóc mái còn rối tung. Cô ngáp dài một cái rồi lơ đãng nói: ``Nếu sau này câu được cá to thì nhớ gọi chị ăn mừng nhé.''

Ryuuhou ngẩng lên hỏi với vẻ tò mò: ``Chị cũng biết cách câu cá à?''

Kson lắc đầu ngay, cười toe toét: ``Không biết câu đâu. Nhưng chị siêu giỏi ăn cá, có bao nhiêu cũng hết.''

Từ trong căn bếp đang bốc mùi thơm của bánh mì nướng, Kuon vọng ra tiếng cười: ``Nếu giỏi ăn cá vậy thì em có thể bắt đầu bằng việc phụ Suzuno chuẩn bị bữa sáng cho cả nhà đi. Khuấy bột, rửa chén gì đó cũng cần người giúp đấy.''

Nàng Tổng trưởng lập tức dừng bước chân định đi vào phòng ngủ, ngoái đầu nhìn về hướng bếp với vẻ mặt hoảng sợ: ``Cái công việc đó nghe khó gấp trăm lần việc học câu cá của Ryuuhou-chan rồi. Thôi em có việc bận, đi đây nhé!''

Tiếng cười lớn của Kaigou vang lên ngay sau khi Kson vội vàng bỏ chạy. Còn Ryuuhou thì đã cúi xuống tập trung đọc từng dòng trong cuốn hướng dẫn, cố gắng thực hiện nút thắt đầu tiên. Nhưng không ngờ sợi dây cước cứ mãi quấn vào ngón tay em, càng cố gắng tháo ra lại càng mắc chặt. Cô bé chỉ còn cách gọi chị mình với giọng bối rối: ``Onee-sama, hình như em làm sai mất rồi. Sao nó cứ mắc vào tay em thế này?''

``Đừng cuộn tròn tay lại nữa, kẻo lại càng mắc chặt thêm,'' cô nàng thở dài bất lực, vừa nói vừa bước đến gần em gái mình. Cô vừa đưa tay ra nhẹ nhàng gỡ từng vòng dây cước ra khỏi các ngón tay đang bị quấn chặt của Ryuuhou, động tác thành thục và cẩn thận để không làm em ấy bị đau.

Ryuuhou rụt rè ngước nhìn chị, mắt to tròn lộ rõ vẻ bối rối khi vừa làm hỏng việc ngay từ đầu: ``Chị\dots\ chị giúp em thắt lại nút này được không? Em làm mãi mà không đúng được một lần nào.''

``Được thôi, chị giúp một lần. Nhưng nhớ cho kỹ nhé, lần sau phải tự làm được, đừng có mãi ỷ lại vào chị nữa,'' Kaigou vừa nói vừa cầm sợi dây cước trên tay, chỉ từng bước một để em gái theo dõi. Những thao tác đơn giản nhưng cần độ chính xác cao được cô thực hiện lướt qua, chỉ trong vài giây đã có một nút thắt chắc chắn, gọn gàng trên đầu dây.

Ryuuhou gật đầu lia lịa, mắt không rời khỏi đôi tay của chị để ghi nhớ từng động tác, hứa sẽ không quên lần sau. Vừa lúc đó, từ trong căn bếp thơm lừng mùi bánh mì nướng và trà nóng, giọng nói trong trẻo của Suzuno vọng ra, gọi cả nhà chuẩn bị vào ăn sáng: ``Mọi người ơi, bữa sáng đã sẵn sàng hết rồi ạ, vào dùng thôi ạ!''

Nghe lời mời, Naomi là người đầu tiên nhảy cẫng lên khỏi ghế. Cô bé nhanh nhẹn nắm lấy tay Koneko, kéo bạn đứng dậy rồi vui vẻ dẫn đường vào phòng ăn. Koneko vẫn ôm chặt quả bóng màu cam trước ngực, vừa đi vừa tíu tít kể cho bạn nghe về những món bánh ngọt mà Suzuno thường làm vào buổi sáng. Tiếng nói cười trong trẻo của hai đứa trẻ khiến bầu không khí trong nhà càng thêm rộn ràng.

Mikeneko cũng thong thả bước theo sau, trên tay cầm chiếc khăn nhỏ vừa dùng để lau những mẩu bánh vụn còn sót lại trên bàn. Thấy Naomi kéo bạn đi quá nhanh, cô khẽ nhắc hai đứa cẩn thận kẻo vấp ngã, nhưng trên môi vẫn nở nụ cười dịu dàng trước vẻ háo hức của chúng. 

Kaigou thu dọn cuộn dây cước rồi cùng Ryuuhou đi vào sau mọi người. Ryuuhou còn cẩn thận kiểm tra lại chiếc hộp dụng cụ câu cá, như muốn chắc chắn rằng nút thắt đầu tiên mình vừa học vẫn được giữ nguyên. Tuy chưa thể tự làm thật thành thạo, cô bé vẫn cảm thấy vui vì đã học thêm được một điều mới.

Chẳng mấy chốc, bàn ăn đã đầy đủ mọi người, tiếng trò chuyện và tiếng cười vang lên hòa cùng mùi bánh mì thơm phức. Chuyến đi câu cá hôm ấy không thành, nhưng trong lòng Ryuuhou đã bắt đầu nhen lên niềm mong đợi về một ngày được tự tay thả cần xuống mặt nước.

Ngày nghỉ rồi cũng sẽ qua, nhưng bộ cần câu vẫn còn đó, nằm gọn trong hộp và chờ đợi một chuyến đi khác, vào một ngày nào đó không xa, khi cả nhà lại cùng nhau tận hưởng những khoảnh khắc vui vẻ bên nhau.

\end{document}